\chapter{真空中的静磁场}
本章我们将考虑一些电场力之外的东西，在第八章的\Refx{定律}{库仑定律}中我们提到，两个电荷间的相互作用力以平方反比减弱，但正如第八章标题静电场所言，这一结果仅适用于静止电荷，当电荷运动时，这一定律就不完全准确了，我们将会看到，运动电荷间的相互作用力会以一种复杂方式依赖于电荷的的速度，而不再仅取决于电荷间的距离。

其中，为了分析上的便利，将运动电荷静止时的那部分力称为电场力，将运动电荷间因运动产生的额外的力称为磁场力，而这些磁场力所遵循的规律，就是本章将研究的内容。

值得注意的是，尽管我们将电荷间的相互作用划分成了电效应和磁效应，但是本章稍后我们就会看到，这两者实际是不可区分的。对于同样的两个电荷，从一些角度看它们间即有电效应也有磁效应，从另一些角度看却只有电效应。因此，电效应和磁效应，电场和磁场，本质是一个统一体的两个方面，哪部分多一些，哪部分少一些，仅取决于观测的角度，而这统一体就是所谓的电磁效应和电磁场。这也是为何这一卷的内容以电磁学为标题的原因。

直接处理电荷间的电磁力是比较复杂的，但是由于一些普遍的原理，因而我们有可能以相对简单的方式来处理电磁力，我们从实验中发现，作用于某一特定电荷上的力，只取决于该电荷的位置、速度、电荷量。我们可以将作用于一个速度为$v$的运动电荷$q$上的力写作
\begin{equation*}
    \vb*{F}=q(\vb*{E}+\vb*{v}\times\vb*{B})
\end{equation*}
上式中的$\vb*{E}$与$\vb*{B}$分别表征电场和磁场，它们是与空间有关的函数。

这样就拆分为出了两个问题
\begin{enumerate}
    \item 电荷在电磁场中受到何种大小和方向的电磁力？
    \item 电荷以何种方式产生怎样的电磁场？
\end{enumerate}
关于这两个问题，有关电场的部分前两章中的静电学中已经很好的回答了，有关磁场的部分将在接下来两章中的静磁学中进行详细的研究，即电荷如何产生磁场并如何受磁场力。

总而言之，\textbf{磁就是电荷因运动产生的附加于电的那一部分效应}。

\section{电流与电流密度}
电荷的运动将在电场之外产生磁场，而大量定向运动的电荷将形成电流，因此在开始本章对磁场的研究前，我们有必要先对电流以及如何定量描述电流分布做一些更多的探讨。

电荷在电场作用下的定向运动形成\uwave{电流}，而形成电流的带电粒子统称为\uwave{载流子}，最为常见的载流子是导体中的电子。但除此之外，在溶液中的载流子还可以是\uwave{阴阳离子}，在半导体中的载流子还可以是带正电的\uwave{空穴}。接下来电流的讨论均以金属导体中的电子为背景展开。

电场作用下，正负电荷所受电场力相反，定向运动的方向也相反，其中我们将正电荷运动的方向规定为电流的方向。例如，虽然在金属导体中的载流子实际是带负电荷的电子，但是我们就可以将等效视为沿反向运动的等量正电荷，此时电流的方向就是电子运动的反方向。\footnote{由此可见，电流并不是载流子的运动方向，而是正载流子的方向或负载流子的反方向。}

由此可见，产生电流的条件是
\begin{enumerate}
    \item 存在可以自由移动的电荷。
    \item 存在电场的作用。
\end{enumerate}
这也就解释了为何电介质是绝缘体，电介质中没有可以自由移动的电荷，也就无法导电了。

现在的问题是，如何定量的描述电流？

\subsection{电流强度}
\uwave{电流强度}是一个我们所熟知的一个用于描绘电流的物理量\footnote{在不引起误会时，电流强度也可以直接简称为电流。}，它反应了电流的强弱，其被定义为单位时间内通过导体某一横截面的电荷量，通常用符号$I$表示，即
\begin{BoxDefinition}[电流强度]
    I=\Fdif{Q}{t}
\end{BoxDefinition}
电流强度在国际单位制中的单位是\si{A=C\cdot s^{-1}}，称为\uwave{安培}，量纲是\si{[I]}。

电流强度是标量，通常所说的电流强度的方向是指电流的流向，即电流是经由界面自左向右流动还是自右向左流动，有时我们也用电流强度的正负来标记电流的流向。

如果我们研究和考虑的情形比较简单，例如电流在一根导线中流动，这根导线的粗细和材质，或许是均匀的，或许在局部有些不均匀，但是如果我们只需要知道单位时间内留过这导线某一截面的电荷量，并且并不关心在那截面上细微而具体的电流分布，那这些（粗细和材质的均匀与否）都不是很重要，电流强度已足矣反映我们关于电流需要知道的一切了。

但是若我们想考虑的更细致一些，如\Refx{图}{电流在粗细不均匀的导线中的流动}所示，很明显，导线较粗处电流也较为稀疏，导线较细处电流就较为密集，但是通过任何一处截面的电流强度是完全相同，那么如果我们想定量的研究一些这种“电流的疏密”，此时电流强度就无能为力了。这尚且还是在导线中的情况，这里电流强度至少还能描述一些东西。

可是在实际中我们还会遇到电流在大块异形导体中流动的情况，例如通过向大地通电流进行地质勘探。此时导体中电流的方向处处都不同，其根本没有一个基本一致的流向，其是完全散开的，这时电流强度就完全无法反映任何信息了，我们甚至都不知道该取哪一截面！

\newpage

\begin{Figure}{电流在粗细不均匀的导线中的流动}
    \inputfigure[scale=0.6]{Chapter10A_Figure01}
\end{Figure}

\subsection{电流密度}
由以上讨论可知，我们有必要在电流强度外定义一个描述电流分布的物理量，这应该是一个空间矢量函数，其方向反映该处电流的流向，其大小反映该处的电流的强弱，很明显，若导体中某一点处的电荷越多且电荷运动速度越快，那么该点处的电流就越强。

由于导体中的电荷是体分布的，因此我们就将导体中某一点处的电荷体密度$\rho$与该点处的电荷运动速度$\vb*{v}$的乘积定义为该点处的\uwave{电流密度}，通常用符号$\vb*{J}$表示，即
\begin{BoxDefinition}[电流密度]
    \vb*{J}=\rho\vb*{v}
\end{BoxDefinition}
电流密度在国际单位制中的单位是\si{A\cdot m^{-2}=C\cdot m^{-2}\cdot s^{-1}}，量纲是\si{[L^{-2}I]}。

通过\Refx{定义}{电流密度}，我们就可以更好理解什么是电流强度了，电流密度描述了电流的分布，电流强度反映的是通过一个截面的总电流量，即\textbf{电流强度是电流密度的通量}。
\begin{BoxEquation}[电流与电流密度]
    I=\Isnt[S]{\vb*{J}\cdot\dif{\vb*{S}}}
\end{BoxEquation}
我们可以用\Refx{图}{电流在粗细不均匀的导线中的流动}中那样的流线图将电流场形象的表现出来，称为\uwave{电流线图}，其疏密反映电流密度的大小，其方向反映电流密度的方向，如果我们在电流线图中取一个截面，那么通过该截面的电流线数目即该截面上的电流强度。

我们现在来看如果所取的“截面”是一个闭合曲面时会发生什么，电流密度是等效正电荷的流动方向，电流此时表示的应是单位时间内流出闭合曲面的电荷量，而根据\Refx{定律}{电荷守恒定律}，电荷不可能凭空出现或消失，因此闭合曲面内的的总电荷量$q$的变化完全来自电荷流入或流出该闭合曲面，故此时\Refx{公式}{电流与电流密度}就可以表示为
\begin{BoxPrinciple}[电流连续性方程]
    流出闭合曲面的电流，即该闭合曲面内的电荷量的负变化率。
    \begin{BoxEq}
        I=\Isot[S]{\vb*{J}\cdot\dif{\vb*{S}}}=-\Fdif{q}{t}
    \end{BoxEq}
\end{BoxPrinciple}
这里的负号是因为电流强度的正值表示电荷向闭合曲面外运动，即为电荷的负变化率，这与先前的\Refx{定义}{电流强度}并不矛盾，那时的电荷量是指流经截面的电荷，意义不同。

这里我们也可以将\Refx{定律}{电流连续性方程}以微分形式表示
\begin{BoxPrinciple}[电流连续性方程的微分形式]
    电流密度的散度，即该处电荷密度的负变化率。
    \begin{BoxEq}
        \nabla\cdot\vb*{J}=-\Fdif{\rho}{t}
    \end{BoxEq}
\end{BoxPrinciple}
这就说明电流的源完全来自电荷的增加与减少
\begin{itemize}
    \item 如果我们在某处看到一个电流的正源，那么该处的电荷一定在不断减少。
    \item 如果我们在某处看到一个电流的负源，那么该处的电荷一定在不断增加。
\end{itemize}\vspace{-15pt}

\subsection{恒定电流场与恒定电场}
现在让我们来考虑一种特殊情况，假设电流密度$\vb*{J}$是与时间无关的，此时电流分布是恒定不变的，称为\uwave{恒定电流}，而我们知道，电流是电荷在电场作用下的定向运动，电流分布不变，电场分布就不变，电荷分布也就不变，即$\rho$和任意闭曲面内的$q$是与时间无关的常数，故
\begin{BoxPrinciple}[电流恒定条件]
    恒定电流场在闭曲面上的通量为零，流入闭曲面的电流总是等于流出闭曲面的电流。
    \begin{BoxEq}
        I=\Isot[S]{\vb*{J}\cdot\dif{\vb*{S}}}=0
    \end{BoxEq}
\end{BoxPrinciple}
\begin{BoxPrinciple}[电流恒定条件的微分形式]
    恒定电流场的散度处处为零，电流分布与电荷分布不随时间变化。
    \begin{BoxEq}
        \nabla\cdot\vb*{J}=0
    \end{BoxEq}
\end{BoxPrinciple}
这就是说，恒定电流场必定是无散的，虽然电荷在导体中在不断流动，但是电荷的分布却没有发生任何变化，从而产生的电场也是恒定不随时间变化的，称之为\uwave{恒定电场}。

恒定电场与静电场具有相似之处，两者的电荷分布和电场分布都不随时间变化，因此描述静电场基本性质的\Refx{定律}{电场高斯定理}和\Refx{定律}{电场环路定理}对于恒定电场也是完全适用的，因此我们也可以在在恒定电场中也可以引入\Refx{定义}{电势}的概念。

恒定电场与静电场最主要的区别在于达成电荷分布不随时间变化的方式，后者的电荷是真正静止的，前者则是通过电荷在定向运动中的动态平衡，即某点处在有电荷流出时必然有等量电荷流入进行补充，从而实现电荷分布不变。因此，静电场中的静电平衡理论，导体内的场强处处为零，导体表面等势，在恒定电场中就不再适用了，在恒定电场中的导体
\begin{itemize}
    \item 导体表面具有一定的电势差，且不随时间变化。
    \item 导体内部具有恒定的电场强度，因而能使电荷定向运动形成恒定电流。
\end{itemize}
恒定电场最直观的一个例子就是电路，导线虽然是一种导体，但在接入电源后，导线却并不是等势的，导线正负两端的电势是不同的。这也反映出了恒定电场的另外一个特征，即恒定电场的维持必定是需要能量的，例如电源。相比之下，静电场的维持不需要能量。

\subsection{恒定电路与欧姆定律}
恒定电流场需要电荷源源不断的流动，这种情形只有在那种首尾相接的路线中才有有可能出现，这种由导线构成的完整回路称为\uwave{电路}，当然，电路并不总是能形成恒定电流场，电路中可以形成恒定电流场的称为\uwave{恒定电路}，电阻电路是恒定电路，电容电路则不是。

% 恒定电路最大的特点是其电学参量是不随时间变化的常量，对于电阻电路，在开关闭合后，电压和电流便是我们计算的那个数字了，现在如此，将来也如此。相反，对于电容电路，正如先前我们推导的\Refx{公式}{电容充电过程的电压}和\Refx{公式}{电容放电过程的电压}所示的那样，电压和电流都是随时间变化的函数，这种非恒定电路的情况就较为复杂了。

恒定电路的一个重要特点是，\textbf{在同一支路上的电流处处相等}，证明是非常容易的，我们可以在导线上取两个截面，连同其间的导线表面取作一个完整的闭合曲面，而后者上显然不会通过任何电流，结合\Refx{定律}{电流恒定条件}可知，从后截面流入的电流总是等于从前截面流出的电流，由于这样的截面是可以任选的，因此任意截面上的电流都是相同的。

% 这一结论是非常重要的，虽然我们使用的导线通常是均匀的，但是导线间连接的电阻器的截面积与导线的截面积必然是不同的，只有通过这一结论，我们才可以坚定的说
% \begin{center}
%     “瞧！电阻上的电流与导线相同！”
% \end{center}
恒定电路中的电流是常数，而实验表明，在大部分情况下，电路两端的电压总是与电流成正比，这一比例系数称为\uwave{电阻}，记作$R$，这一经验结论就是我们非常熟悉的\uwave{欧姆定律}
\begin{BoxPrinciple}[欧姆定律]
    电路两端的电压总是与该电路上的电流成正比。
    \begin{BoxEq}
        U=IR
    \end{BoxEq}
\end{BoxPrinciple}
电阻在国际单位制中的单位是~\si{\ohm=V\cdot A^{-1}}，称为\uwave{欧姆}，量纲是\si{[L^{2}MT^{-3}I^{-2}]}。

电阻的倒数称为\uwave{电导}，记作$G$\vspace{-5pt}
\begin{BoxDefinition}[电导]
    G=\frac{1}{R}
\end{BoxDefinition}
电导在国际单位制中的单位是\si{S=\ohm^{-1}=A\cdot V^{-1}}，称为\uwave{西门子}，量纲是\si{[L^{-2}M^{-1}T^3I^2]}。

应当指出，\Refx{定律}{欧姆定律}并不总是成立的，使得欧姆定律成立的器件称为\uwave{欧姆元件}，否则称为\uwave{非欧姆元件}，其中一个我们熟悉的例子是小灯泡，它的电阻会随着温度变化。\footnote{小灯泡本质仍然是电阻丝，只不过为了发光，在电阻丝上的温度变化范围远超正常的电阻器，以至于电阻率有了明显变化。}

导体电阻的大小与导体的材料和几何形状有关，对于由一定材料制成的横截面均匀的导体，电阻与导体长度$L$成正比，电阻与导体横截面积$S$成反比，这一结论即\uwave{电阻定律}
\begin{BoxEquation}[电阻定律]
    R=\rho\frac{L}{S}
\end{BoxEquation}
上式中的$\rho$不是电荷密度，其表示材料的\uwave{电阻率}，其倒数称为\uwave{电导率}。
\begin{BoxDefinition}[电导率]
    \gamma=\frac{1}{\rho}
\end{BoxDefinition}
电阻率在国际单位制中的单位是\si{\ohm\cdot m^{-1}}，量纲是\si{[LMT^{-3}I^{-2}]}。

电导率在国际单位制中的单位是\si{S\cdot m}，量纲是\si{[L^{-1}M^{-1}T^{-3}I^{-2}]}。

现在让我们来研究一下\Refx{定律}{欧姆定律}在微观上会导致什么样的结果，我们在导体中取一个很小的\uwave{电流管}，即沿电场线的一个管状假象曲面，由于这里取的电流管很小，可以近似视为一个柱体，设电流管的长度为$\dif{L}$，设电流管的截面积为$\dif{S}$，那么
\begin{Equation}[欧姆定律的微分形式1]
    \dif{I}\xlongequal{\text{\Refx{定律}{欧姆定律}}}\frac{\dif{U}}{R}
    \xlongequal[\text{\Refx{定义}{电导率}}]{\text{\Refx{公式}{电阻定律}}}
    \gamma\dif{U}\frac{\dif{S}}{\dif{L}}
\end{Equation}
其中$\dif{I},\dif{U}$分别是小电流管上的电流和其两端的电压。

由于小电流管中的电场可以近似视为匀强电场，因而
\begin{Equation}[欧姆定律的微分形式2]
    \dif{U}=E\dif{L}
\end{Equation}
将式\ref{式:欧姆定律的微分形式2}代入\ref{式:欧姆定律的微分形式1}，消去$\dif{L}$
\begin{Equation}[欧姆定律的微分形式3]
    \dif{I}=\gamma E\dif{S}
\end{Equation}
由于小电流管截面上的电流可以近似认为是均匀分布，故有$\dif{I}=J\dif{S}$，代入上式约去$\dif{S}$
\begin{equation*}
    J=\gamma E
\end{equation*}
考虑到电流密度的方向应与场强保持一致
\begin{BoxGeneral}[欧姆定律的微分形式][定律]
    \begin{BoxEq}
        \vb*{J}=\gamma\vb*{E}
    \end{BoxEq}
\end{BoxGeneral}
值得注意的是，上式虽然是由恒定电路的情形导出，但是对于那些非恒定电路，由于\Refx{定律}{欧姆定律}在电阻两端仍然是成立的，上式实际上也是适用的\footnote{只不过此时$\vb*{J},\vb*{E}$都会随时间变化，但是两者间的正比关系$\vb*{J}=\gamma\vb*{E}$始终成立，就像$U=IR$在$U,I$变化时也成立。}。该结论在电磁场中的地位与上一章中\Refx{公式}{电位移矢量与场强的电容率表示}是相当的，但同样，这也是对物质特性的一种近似，对于那些非欧姆元件，这一结论就未必正确了。我们仍然以小灯泡为例来说明，如果认为决定电阻的温度只与电流有关，那么此时$\vb*{J},\vb*{E}$是非线性的，如果我们还要考虑温度在小灯泡上的积蓄，那么此时$\vb*{J},\vb*{E}$之间就不存在简单的单值对应关系了。
\section{磁场与磁感强度}
在本章开始时，我们用了相当大的篇幅来强调，磁是运动电荷在电之外的那一部分效应，但是这样的磁似乎与我们熟悉的磁铁的那种磁并不相符，这实际上是一个历史问题。 

在历史上，有关磁现象的研究要远早于电现象，在公元前6世纪我国古代先民就已经懂得用磁石吸引铁器，在公元18世纪后西方科学家才开始逐渐研究电学现象。而磁现象的发现正是始于磁铁的磁性，我们观察到无论是天然磁铁还是人造磁铁，都具有吸引铁、镍、钴及其合金的性质，这种性质称为\uwave{磁性}，这些能被磁铁吸引的物质则统称为\uwave{铁磁质}。

我们发现，条形磁铁两端磁性特别强，称为\uwave{磁极}，条形磁铁中间磁性很弱，称为\uwave{中性区}，另外一个有趣的现象是，如果我们将条形磁铁悬挂起来，使它能在水平面内自由转动，我们观察到，无论我们条形磁铁的最初指向是什么，在稳定后两端的磁极总是分别指向南北方向，并且，总是一个磁极指向南方，而另外一个磁极指向北方，这就说明两个磁极是有差异的。

为了便于区分磁铁的磁极，将磁铁指北的一端称为\uwave{北极}，将磁铁指南的一端称为\uwave{南极}，两者分别可以用字母N和S表示，基于这样的磁极定义，我们还发现如果我们尝试将两根磁铁相互靠近，我们将会看到这样的现象，\textbf{同号磁极相互排斥，异号磁极相互吸引}，这一实验现象说明，地球本身就是一个巨大的磁铁，北极是它的磁南极，南极是它的磁北极。不过实际上地球的地理南北极与磁南北极并不是完全重合的，两者间的夹角称为\uwave{地磁偏角}。

在1785年库仑提出关于\uwave{电的库仑定律}，即我们熟悉的\Refx{定律}{库仑定律}的同时，也提出了\uwave{磁的库仑定律}，它认为磁现象是由于\uwave{磁荷}产生的，就像电现象是由于电荷产生的
\begin{itemize}
    \item 电荷将会产生电场，电荷有“正负”两种，同号电荷相斥，异号电荷相吸，电荷间的相互作用力，与两者电荷量的乘积成正比，与电荷间距离的平方成反比。
    \item 磁荷将会产生磁场，磁荷有“北南”两种，同号磁荷相斥，异号磁荷相吸，磁荷间的相互作用力，与两者磁荷量的乘积成正比，与磁荷间距离的平方成反比。
\end{itemize}
在磁荷磁的库仑定律的帮助下，有关磁铁的现象基本可以被定性解释和定量计算。

\begin{Figure}{磁的磁荷观点}
    \subfigure[磁化前]
    {
        \inputfigure[scale=0.45]{Chapter10A_Figure03}
    }\\
    \subfigure[磁化后]
    {
        \inputfigure[scale=0.45]{Chapter10A_Figure02}
    }\\
    \subfigure[磁化的宏观效果]
    {
        \inputfigure[scale=0.45]{Chapter10A_Figure04}
    }
\end{Figure}

磁荷理论中，磁铁当然不是由两个巨大的“北磁荷”和“南磁荷”拼接而成的，且不深究两者是如何克服巨大斥力结合在一起的，我们知道，磁铁的一个非常奇特的性质就是，如果我们将磁铁沿中性区切成两块，那么我们不会得到两个独立的南北极，而是产生了两个具有完整南北极的小磁铁。因此实际上，磁铁在磁荷理论中被认为是这样一种物质，其内部存在有大量排列整齐的微小磁铁，这种微小磁铁也称为\textbf{磁偶极子}，这类似于电场中的\textbf{电偶极子}。

磁铁的中间部分，磁偶极子的南极和北极总是相互抵消，从而对外在宏观上不显出磁性，这就是所谓的中性区，磁铁两端处位于端面附近的那些磁偶极子，具有大量无法抵消的北磁荷和南磁荷，从而就形成了具有强烈的磁性的北磁极和南磁极，就像被极化的电介质那样。

磁铁的任何一个磁极都可以吸引铁磁质，而铁磁质本身没有任何磁性。这就是因为铁磁质内部同样存在大量的磁偶极子，但是这些磁偶极子原先排列非常混乱，从而在宏观上不显示出任何磁性。而在磁铁的外加磁场作用下，铁磁质中的磁偶极子排列变得整齐，铁磁质就被磁化了，这个过程与电介质中的电偶极子在电极化的过程中排列整齐是十分相似的。但是这里的铁磁质作为一种磁介质按照原先电介质的观点看是较为特殊的，因为在磁化后，即便外磁场褪去，这些\textbf{铁磁质}也往往可以保留一定磁性，这相当于电介质中的\textbf{铁电体}。有一部分铁磁质时间长了或是经历碰撞和加热之后就会被消磁，有一部分铁磁质则可以更为稳定的长久保留磁性，这类磁介质也称为\textbf{永磁体}，这相当于电介质中的\textbf{永驻体}，实际上，我们所使用的天然磁铁就是一种天然的永磁体，而所谓的人造磁铁就算通过磁化的方式制造的永磁体。

电学曾定义了电场强度$\vb*{E}$，磁学中也可以相仿的定义磁场强度$\vb*{H}$，即单位磁荷在磁场中所受到的力（我们或许需要将北磁荷和南磁荷分别赋予正和负），通过点磁荷的库仑定律可以先推导出点磁荷的磁场，进而计算出磁偶极子的磁场，这样就可以定量分析磁铁了。

\textbf{然而，如果以上这一切都是真的，磁学就会比今天所看到的简单太多了。}

磁荷理论的一个缺陷在于，\uwave{磁单极子}从未被发现，即独立的北磁荷和南磁荷并不存在，如果我们将一个磁铁切开只会得到更多的小磁铁，也没有任何一种粒子具有单独的磁荷。而相比之下在电学中，\uwave{电单极子}则是常见的，即独立的正电荷和负电荷，质子和电子就是最容易看到的例子。然而实际上，磁单极子不存在并非是磁荷理论被历史舍弃的主要原因，因为我们完全可以退一步认为所有磁荷都是以类似于束缚电荷那样的束缚磁荷的形态存在的。

磁荷理论的真正问题在于，电流和运动电荷的磁效应的发现。

在1820年，丹麦科学家奥斯特通过著名的\uwave{奥斯特实验}，发现将一个可以自由旋转的小磁针置于直导线的正下方，然后给直导线通以电流，观察到小磁针将发生偏转，而且当电流方向改变时，观察到小磁针的偏转方向将相应的发生反转，\textbf{这就说明了电流同样可以产生磁场}。

在这里有一个非常相似的实验，如果将一个根水平导悬挂在一个马蹄形磁铁的磁极之间，然后给直导线通以电流，观察到水平导线将会发生偏移，\textbf{这就说明了电流会受到磁场作用}。

在这个跨时代的实验的消息抵达法国科学院的仅1周之后，法国科学家安培通过实验，发现两根平行载流导线之间也存在相互作用力，如果导线的电流方向相同则相互吸引，如果导线的电流方向相反则相互排斥，\textbf{这就说明电流之间可以独立的通过磁场产生相互作用}。

电流是电荷定向运动形成的，通过实验我们也可以观察到磁场对运动电荷的作用，阴极射线管是一个真空放电管，在它两个电极之间加上高电压时，在它的阴极就会产生所谓的阴极射线，也就是电子束。电子束本身是不可见的，为此在阴极射线管中附有可以观测电子束轨迹的荧光屏，通常情况下电子束是沿直线前进的，但是如果我们在阴极射线管旁放置一块磁铁，就会观察到电子束的轨迹发生了偏转，这就可以表明电子束的受到了磁场作用。

磁铁和磁铁间，电流和磁铁间，电流和电流间，都存在相互作用，当然或许这其实是三种不同的作用，但是从自然规律的统一性上看我们更愿意相信这实际是同一种磁相互作用，然而这就导致了一个问题，尽管我们可以称磁铁中有一些磁偶极子，但是电流中绝无可能有任何的磁偶极子，也就是说，磁荷理论无法解释电流的磁效应，至此，磁荷理论也就被抛弃了。

在今天，我们从一个相反的观点来看待磁现象，我们认为磁场的本源是电流或者说是电流中的运动电荷，这也就是本章开始处“磁是运动电荷在电之外的那一部分效应”的含义，那么在这种新的认识下，磁铁的磁效应又应该如何解释呢？安培提出了分子电流假说，这一假说认为，磁铁在分子层面中有许多排列整齐的微小环状电流，这些取向一致分子电流在宏观上等效于一个环绕磁铁边缘的大环状电流，从而对外显示出磁性，磁铁的磁化实际上就是将原先铁磁质中排列混乱的分子电流通过磁场作用排列整齐，从而使得铁磁质也具有磁性。
\begin{Figure}{磁的分子电流观点}
    \subfigure[磁化前]
    {
        \inputfigure[scale=0.45]{Chapter10A_Figure06}
    }\\
    \subfigure[磁化后]
    {
        \inputfigure[scale=0.45]{Chapter10A_Figure05}
    }\\
    \subfigure[磁化的宏观效果]
    {
        \inputfigure[scale=0.45]{Chapter10A_Figure07}
    }
\end{Figure}

安培分子电流假说与实验事实是相符的，因为依照它的观点，磁铁像是一个由分子电流构成的通电螺旋管，而实际上，磁铁产生的磁场与真正的通电螺旋管的磁场是非常相似的。

安培分子电流假说不仅与电流的磁效应兼容，而且其与实际的物质微观结构是相符的，我们现在知道，电子在环绕原子核转动的同时，电子本身也会发生自旋，这两者实际就是分子电流的微观来源。当然这种“电子绕核转动”和“电子自旋”只是经典观点下提法，现在我们已经知道，电子实际是以电子云的形态存在，电子自旋也只是一个粒子属性，并非是像真正的宏观物体那样的自转，这里最为本质的研究需要用到量子力学，但是经典力学的这些提法具有启发性的，尽管在原理上它不完全正确，但足矣对物质磁性给出足够准确的猜测了。

本章将首先研究真空中的磁场，而有关那些磁介质和磁铁以及安培分子电流假说的内容将在下一章讨论。暂时我们只需要理解，即便现在我们认为磁场是源于电流，磁铁的磁场仍然是同样客观真实的，磁铁的磁效应最终也可以用电流解释，并非什么理论之外的奇怪东西。

\subsection{磁感应强度}
在电场中，我们引入了描述静止电荷间的电相互作用力的\Refx{定律}{库仑定律}，随后通过该结论推导出试探电荷在空间中某一特定位置受到的电场力仅与其电荷量有关，即
\begin{Equation}[电场力的作用]
    \vb*{F}=q\vb*{E}
\end{Equation}
在磁场中，由于一些原因，我们很难通过实验测定运动电荷间的磁相互作用力的公式，但是我们发现，如果现在有一些磁场，或许是由旁边的电流和磁铁产生的，这不重要，那么通过实验将观察到试探电荷在空间某一特定位置受到的磁场力仅与其电荷量和速度有关，即\footnote{这里要假定空间中没有任何电场，或者我们能计算出电场力的大小并在试探电荷的受力中减去，因为实验中我们并不能分辨试探电荷受到的作用力中那一部分属于电场力那一部分属于磁场力。这里研究的是磁场，我们只需要其中磁场力的那一部分。}
\begin{Equation}[磁场的作用]
    \vb*{F}=q\vb*{v}\times\vb*{B}
\end{Equation}
式\ref{式:磁场的作用}要比式\ref{式:电场力的作用}复杂，电场力仅和电荷量有关，磁场力同时与电荷量和速度有关。

式\ref{式:磁场的作用}中的矢量$\vb*{B}$与先前电场中的电场强度$\vb*{E}$的地位是相似的，其定量反映了磁场的空间分布，称为\uwave{磁感应强度}或简称为\uwave{磁感强度}。磁感强度的确定可以通过测量试探电荷在磁场中的磁场力实现，令试探电荷$q$以某个大小相同的速度$\vb*{v}$经过同一位置，将试探电荷$q$受到的最大磁场力记为$F_{\max}$，根据叉乘原理，此时$\vb*{F}_{\max},\vb*{v},\vb*{B}$垂直且满足$F_{\max}=qvB$，故
\begin{BoxDefinition}[磁感强度]
    \vb*{B}=\frac{\vb*{F}_{\max}\times\vu*{v}}{qv}
\end{BoxDefinition}
这里的矢量关系可以由下图看出，即当磁场力最大$\vb*{F}=\vb*{F}_{\max}$时，此时磁感强度$\vb*{B}$垂直于由$\vb*{F},\vb*{v}$构成的平面，且磁感应强度$\vb*{B}$的方向为右手四指由$\vb*{F}$转向$\vb*{v}$时大拇指的方向。
\begin{Figure}{磁感强度的矢量关系}
    \subfigure[当$\vb*{v},\vb*{B}$为垂直时的$\vb*{F}$视角]
    {
        \hspace{30pt}
        \inputfigure{Chapter10B_Figure01}
        \hspace{30pt}
    }
    \subfigure[当$\vb*{v},\vb*{B}$为垂直时的$\vb*{B}$视角]
    {
        \hspace{30pt}
        \inputfigure{Chapter10B_Figure02}
        \hspace{30pt}
    }\\
    \subfigure[当$\vb*{v},\vb*{B}$不垂直时的$\vb*{F}$视角]
    {
        \hspace{30pt}
        \inputfigure{Chapter10B_Figure03}
        \hspace{30pt}
    }
    \subfigure[当$\vb*{v},\vb*{B}$不垂直时的$\vb*{B}$视角]
    {
        \hspace{30pt}
        \inputfigure{Chapter10B_Figure04}
        \hspace{30pt}
    }
\end{Figure}

这里我们也可以看到，在一般情况下，实际上$\vb*{B}$并不总是垂直于$\vb*{F},\vb*{v}$构成的平面，因为先验成立的是$\vb*{F}=q\vb*{v}\times\vb*{B}$，所以只有当三者垂直磁场力最大时我们才能逆向求出磁感强度。

磁感强度在国际单位制中的单位是\si{T=N\cdot A^{-1}\cdot m^{-1}}，称为\uwave{特斯拉}，量纲是\si{[MT^{-2}I^{-1}]}。

\textbf{磁感强度}描述了磁场分布，虽然其地位等同于电场强度，但是不能将其称为\textbf{磁场强度}，这就是因为本节初叙述的历史原因，磁场强度$\vb*{H}$的名称已经在磁荷理论中被先行被定义为了单位磁荷在磁场中受到的力，因而后来基于试探运动电荷的定义只能改称为磁感强度$\vb*{B}$。

\subsection{磁感强度的叠加原理}
实验表明，当空间中有若干个磁场源的情况下，它们产生的磁场服从叠加原理，它们的总磁感强度就是各个磁场源独立存在时的磁感强度$\vb*{B}_1,\vb*{B}_2,\cdots,\vb*{B}_n$的矢量和，即
\begin{BoxPrinciple}[磁感强度的叠加原理]
    若干磁场源在磁场中某点处产生的磁感强度，

    等个各磁场源单独存在时，在该点处产生的磁感强度的矢量和。
    \begin{BoxEq}
        \vb*{B}=\SumIN{\vb*{B}_i}
    \end{BoxEq}
\end{BoxPrinciple}
\Refx{定律}{磁感强度的叠加原理}与先前\Refx{定律}{电场强度的叠加原理}是电磁场中最基本的两个规律，电场中的这一点是我们从力的叠加原理推证出来的，磁场中如果我们愿意的话其实仍然可以这样做，但实际上，场的叠加原理才是更为普遍和基本的规律。\footnote{场的叠加原理是更普遍和基本的，因此从力的叠加原理推导它实在没有什么必要，唯一的价值是说明两者不矛盾。}

\section{毕奥-萨伐尔定律}

\subsection{毕奥-萨伐尔定律}
本节我们来研究运动电荷或运动电荷构成的电流是如何产生磁场的，然而由于一些后面我们会提到的原因，运动电荷产生的磁场很难被实验观测到，运动电荷构成的电流又往往是具有特定形状的，例如通电直导线或通电线圈，无法归纳出普适性比较强的规律。

在安培发现两根载流导线之间的磁相互作用力后，另外两位法国科学家毕奥和萨伐尔随后提出了定量描述电流磁场的\uwave{毕奥-萨伐尔定律}。这里引入了\uwave{电流元}的概念，我们知道，由于磁场适用叠加原理，我们就可以将任意形状的电流视为无穷多段无限小的电流的集合，这些无限小的电流就是所谓的电流元，因此只要能弄清电流元产生的磁场（这就是毕奥-萨伐尔定律的具体内容），那么任意形状的电流的磁场就可以通过电流元的磁场的叠加求得了。

我们需要指出的是，电流元虽然长度无限小，但是仍然是有方向的，故通常记作$I\dif\vb*{l}$。

我们现在来看毕奥-萨伐尔定律是如何表述的
\begin{BoxGeneral}[毕奥-萨伐尔定律][定律]
    \begin{BoxEq}
        \dif{\vb*{B}}=\frac{\mu_0}{4\pi}
        \frac{I\dif{\vb*{l}}\times\vu*{r}}{r^2}
    \end{BoxEq}
\end{BoxGeneral}
\Refx{定律}{毕奥-萨伐尔定律}指出，电流元产生的磁场垂直于$I\dif{\vb*{l}},\vb*{r}$构成的平面，且指向右手四指由$I\dif{\vb*{l}}$转向$\vb*{r}$时大拇指的方向，其大小，与距离平方成反比，与电流成正比。

\Refx{定律}{毕奥-萨伐尔定律}中$\mu_0$为真空磁导率，其数值为$\mu_0=4\pi\times 10^{-7}$\si{N\cdot A^{-2}}，这里真空磁导率的数值之所以如此齐整，是因为和安培的定义有关，稍后我们会进一步讨论。

\Refx{定律}{毕奥-萨伐尔定律}并不总是成立的，其仅在\uwave{静磁场}的情况下精确成立，这实际上要求产生磁场的电流是一个恒定电流场，即组成电流的每个电流元的大小和方向不随时间变化，这与先前电荷间的\Refx{定律}{库仑定律}只能适用于\uwave{静电场}的情况是类似的。

\Refx{定律}{毕奥-萨伐尔定律}可以根据磁场叠加原理计算任意载流导线的磁场
\begin{BoxEquation}[电流的磁场]
    \vb*{B}=\Ilnt[L]{\dif\vb*{B}}=\Ilnt[L]{\frac{\mu_0}{4\pi}
    \frac{I\dif{\vb*{l}}\times\vu*{r}}{r^2}}
\end{BoxEquation}
这里只需要曲线积分，因为目前我们考虑的电流都是在导线中的，即电流都是线状的。

值得注意的是，\Refx{定律}{毕奥-萨伐尔定律}中设想的电流元只是一种微元模型，电流元是不可能单独存在的，因此该定律是无法直接通过物理实验验证的，但是通过其积分得到的各类不同载流导线的磁场都与实验都相符，这就间接验证了定律的正确性。

\subsubsection{载流直导线的磁场}
设有一载流直导线，长度为$L$，电流为$I$，现考察载流直导线外的场点$P$处的磁感强度。

设$P$点到导线的垂足$O$为原点，以直导线上的电流方向为$x$轴的正方向建立平面直角坐标系，在该坐标系中点$P$在$y$轴上，且设其离开导线的垂直距离（即纵坐标）为$y$。

设$P$点和导线两端的连线与导线之间的夹角分别为$\theta_1$和$\theta_2$。
\begin{Figure}{载流直导线的磁场}
    \inputfigure[scale=0.85]{Chapter10B_Figure05}
\end{Figure}
我们可以将载流直导线视为一系列的电流元的组合，根据\Refx{定律}{毕奥-萨伐尔定律}可知，导线上方的磁场垂直直面向外，导线下方的磁场垂直直面向内，即在同一场点处，每个电流元产生的磁场方向都相同，因此我们就可以将矢量$\dif{\vb*{B}}$转化为标量$\dif{B}$计算。

这里为了方便在平面图形上表示那些垂直于纸面的方向，如\Refx{图}{载流直导线的磁场}所示，我们往往会用$\odot$和$\otimes$表示\uwave{垂直纸面向外}和\uwave{垂直纸面向内}的概念\footnote{实际上这种记号在磁场更多是用于表达垂直直面的电流，而不是磁感强度。}，这种记号可以通过离弦的箭来记忆，射向自己的箭将看到点状箭尖，射出的箭将看到十字形的尾部箭羽。

这里以载流直导线上方的场点$P$为例，将\Refx{定律}{毕奥-萨伐尔定律}改写为标量形式
\begin{Equation}[载流直导线1]
    \dif{B}=\frac{\mu_0}{4\pi}\frac{I\dif{l}\sin\theta}{r^2}
\end{Equation}
式\ref{式:载流直导线1}无法直接积分，因为其关于$\theta$但却具有微分$\dif{l}$，因此要首先将$\dif{l}$转换为$\dif{\theta}$。

关注到$l=r|\cos\theta|$和$y=r|\sin\theta|$，因此
\begin{equation*}
    \frac{l^2}{y^2}=\cot^2\theta\qquad\frac{l}{y}=\abs{\cot\theta}
\end{equation*}
因此$r^2$可以被表示为
\begin{Equation}[载流直导线2]
    r^2=y^2+l^2=y^2(1+\cot^2\theta)=y^2\csc^2\theta
\end{Equation}
因此$l$的微分$\dif{l}$可以被表示为
\begin{Equation}[载流直导线3]
    l=y|\cot\theta|\qquad \dif{l}=y|\csc^2\theta\dif{\theta}|=y\csc^2\theta\dif{\theta}
\end{Equation}
这就有
\begin{Equation}[载流直导线4]
    \frac{I\dif{l}}{r^2}=\frac{Iy\csc^2{\theta}}{y^2\csc^2\theta}=\frac{I\dif{\theta}}{y}
\end{Equation}
式\ref{式:载流直导线4}代入式\ref{式:载流直导线1}中，得到
\begin{Equation}[载流直导线5]
    \dif{B}=\frac{\mu_0}{4\pi}\frac{I\sin{\theta}\dif{\theta}}{y}
\end{Equation}
这就将自变量$\theta$和微分$\dif{\theta}$相统一。

对$\dif{B}$积分，得到
\begin{equation*}
    B=\Int[\theta_1][\theta_2]{\dif{B}}=
    \frac{\mu_0}{4\pi}\frac{I}{y}
    \Int[\theta_1][\theta_2]{\sin{\theta}\dif{\theta}}=
    \frac{\mu_0}{4\pi}\frac{I}{y}(\cos\theta_2-\cos\theta_1)
\end{equation*}
即
\begin{BoxEquation}[载流直导线的磁场]
    B=\frac{\mu_0}{4\pi}\frac{I}{y}(\cos\theta_2-\cos\theta_1)
\end{BoxEquation}
这里，如果载流直导线变为无限长，那么$\theta_1\rightarrow 0,\theta_2\rightarrow\pi$
\begin{BoxEquation}[载流无限长直导线的磁场]
    B=\frac{\mu_0}{2\pi}\frac{I}{y}
\end{BoxEquation}
这里我们可以通过\Refx{图}{载流直导线的磁场}将载流直导线附近的磁场方向归纳如下，即将右手大拇指指向电流方向，那么右手四指的环绕方向即为磁场方向，称为\uwave{右手螺旋法则}。

\subsubsection{载流圆线圈的磁场}
设有一载流圆线圈，半径为$R$，电流为$I$，现考察其轴线上的场点$P$处的磁感强度。
\begin{Figure}{载流圆线圈的磁场}
    \inputfigure[scale=0.64]{Chapter10B_Figure06}
\end{Figure}
设线圈的圆心$O$为原点，以轴线为$x$轴，以线圈所在平面为$yz$平面，建立空间直角坐标系，场点$P$至$yz$平面的距离为$x$，场点$P$至线圈任意一点的距离均为$r=\sqrt{R^2+x^2}$。

这里$\vb*{I}\dif{\vb*{l}},\vb*{r}$总是垂直的，因而根据\Refx{定律}{毕奥-萨伐尔定律}
\begin{equation*}
    \dif{B}=\mk\frac{I\dif{l}}{r^2}
\end{equation*}

如\Refx{图}{载流圆线圈的磁场}所示，由于线圈的对称性，在线圈各处的电流元对轴线上场点施加的磁感强度中，垂直$x$轴的各分量相互抵消，只有平行$x$轴的分量才有贡献。

由于$\cos\alpha=\dfrac{R}{r}$，且$r=\sqrt{R^2+x^2}$
\begin{equation*}
    \dif{B}_x=\dif{B}\cos\alpha=\mk\frac{I\dif{l}}{r^2}\frac{R}{r}=\mk\frac{IR\dif{l}}{r^3}=\mk\frac{IR\dif{l}}{(R^2+x^2)^{\frac{3}{2}}}
\end{equation*}
从而积分得
\begin{equation*}
    B=B_x=\mk\frac{IR\dif{l}}{\qty(R^2+x^2)^{\frac{3}{2}}}\Int[0][2\pi R]{\dif{l}}=
    \frac{\mu_0}{2}\frac{IR^2}{\qty(R^2+x^2)^{\frac{3}{2}}}
\end{equation*}
即
\begin{BoxEquation}[载流圆线圈的磁场]
    B=\frac{\mu_0}{2}\frac{IR^2}{\qty(R^2+x^2)^{\frac{3}{2}}}
\end{BoxEquation}
对于场点位于线圈圆心的情况，这里有$x=0$，即
\begin{BoxEquation}[载流圆线圈在圆心处的磁场]
    B=\frac{\mu_0}{2}\frac{I}{R}
\end{BoxEquation}
对于场点远离线圈的情况，这时有$x\gg R$，即
\begin{BoxEquation}[载流圆线圈在远场点sz的磁场]
    B=\frac{\mu_0}{2}\frac{IR^2}{x^3}
\end{BoxEquation}
以上我们是以场点$P$位于线圈上方的情况考虑的，但是通过\Refx{图}{载流圆线圈的磁场}的下半部分可以看出，这样的推导对线圈下方的场点也是成立的，而且无论场点在线圈上方还是线圈下方，其磁场方向都是沿着线圈电流的右手螺旋方向，这就为矢量表达创造了机会。

记线圈的面积为$S=\pi r^2$，定义载流线圈的\uwave{磁偶极矩}为
\begin{BoxDefinition}[磁偶极矩]
    \vb*{m}=IS\vb*{e_n}
\end{BoxDefinition}
其中$\vb*{e}_n$为载流线圈电流的右手螺旋方向上的单位矢量。
\begin{Figure}{载流圆线圈与磁偶极子}
    \inputfigure[scale=0.64]{Chapter10B_Figure07}
\end{Figure}

那么对于载流线圈的远场点，磁感强度就可以表示为
\begin{BoxEquation}[磁偶极子在中轴线上的磁场]
    \vb*{B}=\frac{\mu_0}{2\pi}\frac{\vb*{m}}{x^3}
\end{BoxEquation}
\Refx{公式}{磁偶极子在中轴线上的磁场}与先前\Refx{公式}{电偶极子延长线上的电场}在形式上很接近，出于这样的原因，仿照电偶极子，我们也将载流线圈称为\uwave{磁偶极子}，因此我们在这里的\Refx{定义}{磁偶极矩}在地位上就与先前的\Refx{定义}{电偶极矩}是相当的了。

磁偶极子和电偶极子在两者偶极矩方向上的远场点处的场的数学形式非常相似，这一点已经通过数学推导证明了，但在下一节我们还会通过场线图的对比看到，磁偶极子和电偶极子在任意远场点处的场的形态也都非常相似。我们或许会感到困惑，因为虽然电场和磁场都是平方反比场，但是两者的性质完全不同，电场是有源无旋场，磁场是有旋无源场，为何两种性质截然不同的场在这里可以产生相同的数学形式？这实际就和偶极子中远场点的要求是密不可分的，因为当距离很远时，无论是源还是旋都会归于零，从而消除了差异。

磁偶极子可以被称为“偶极子”的另一个原因是，载流线圈与先前磁荷理论中同样称为磁偶极子的小磁铁产生的磁场是相当的，即，\textbf{载流线圈的磁场与小磁铁的磁场是等效的}，这再次体现了磁荷理论与分子电流理论在解释磁铁的现象时的一致性，磁铁的磁性总是来自于材料中微小的磁偶极子，无论你认为磁偶极子到底是微小的载流线圈还是微小的磁铁。\footnote{载流线圈的磁场反映的其实就是环状电流的磁场，只不过微观上的分子环状电流其实已经不再需要线圈作为载体了。}
\begin{Figure}{载流线圈与小磁铁的磁场对比}
    \subfigure[载流线圈的磁场]
    {
        \includegraphics[scale=0.8]{old/Python/build/VFP_B03.pdf}
    }\qquad
    \subfigure[小磁铁的磁场]
    {
        \includegraphics[scale=0.8]{old/Python/build/VFP_M03.pdf}
    }
\end{Figure}
\vspace{-20pt}

\subsubsection{载流螺线管的磁场}
螺线管是指均匀地密绕在直圆柱面上的螺旋形线圈，设螺线管的半径为$R$，设螺线管中通有电流$I$，沿轴向单位长度的线圈数为$n$，现在来考察其管内轴线上一点的磁场。

螺线管在一般情况下的磁场是较为复杂的，为了便于计算，我们做几点假定。

我们不难想象，对于真实的螺线管，导线每环绕圆柱面一周，导线在轴线方向就要前进一些距离，这在数学上处理起来就很复杂了，但是如果螺线管上的导线缠绕的非常紧密，导线在同一周上的轴线位移可以忽略，这就使得我们可以将完整的螺旋线视为由一系列分离的圆线圈构成，这就是为何我们将螺线管的剖视图绘制为\Refx{图}{载流螺线管的磁场}的原因。

我们还要假定的一点时，螺旋管上的导线缠绕的如此之紧密，以至于即便我们沿着轴线方向取很小的一段距离，在其上的线圈数目也是十分多的。在这里我们实质是将螺线管上的分立线圈数目通过物理无穷小的想法连续化了，目的是为了使得我们可以通过积分方法叠加各个线圈产生的磁场。当然，这会忽略在导线直径尺度上磁场的不连续性，不过现在的螺线管的线径往往是很细的，在几厘米的距离上就可以缠绕数千圈，这种误差便也可以接受了。

\begin{Figure}{载流螺线管的磁场}
    \inputfigure[scale=0.8]{Chapter10B_Figure08}\vspace{-15pt}
\end{Figure}

以$P$为原点，以螺线管电流的右手螺旋方向为$x$轴，建立平面直角坐标系。现在我们在螺线管上距离$P$点为$l$处的位置取一小段$\dif{l}$上的线圈，在这$\dif{l}$上共有线圈$n\dif{l}$匝，那么根据\Refx{公式}{载流圆线圈的磁场}，在这$\dif{l}$上的线圈在场点$P$处产生的磁感强度$\dif{B}$为
\begin{equation*}
    \dif{B}=\frac{\mu_0}{2}\frac{InR^2\dif{l}}{\qty(R^2+l^2)^{\frac{3}{2}}}
\end{equation*}

在螺线管上各个小段$\dif{l}$上的线圈在场点$P$处产生的磁感强度$\dif{B}$方向均相同，指向电流的右手螺旋方向，因此整个螺线管的磁场即为（设$l_1,l_2$是以$P$为中心螺线管两端的横坐标）
\begin{Equation}[载流螺线管1]
    B=\Int[l_1][l_2]{\frac{\mu_0}{2}\frac{InR^2\dif{l}}{\qty(R^2+l^2)^{\frac{3}{2}}}}
\end{Equation}
式\ref{式:载流螺线管1}关于$\dif{l}$的积分较难积出，为此我们使用三角换元，假设位于$l$处的螺线管相对于场点$P$处的角度为$\beta$，根据螺线管中的几何关系可以看到
\begin{equation*}
    l=R\cot\beta\qquad \dif{l}=-R\csc^2\beta\dif{\beta}
\end{equation*}

另外一方面
\begin{equation*}
    l^2+R^2=R^2\csc^2\beta
\end{equation*}

这样就有
\begin{Equation}[载流螺线管2]
    \frac{InR^2\dif{l}}{\qty(R^2+l^2)^{\frac{3}{2}}}
    =-\frac{InR^3\csc^2\beta\dif{\beta}}{R^3\csc^3\beta}
    =-\frac{In\dif{\beta}}{\csc\beta}
    =-In\sin\beta\dif{\beta}
\end{Equation}
假设$l_1,l_2$对应的角为$\beta_1,\beta_2$，那么
\begin{Equation}[载流螺线管3]
    B=\frac{\mu_0}{2}In\Int[\beta_1][\beta_2]
    {-\sin{\beta}\dif{\beta}}
\end{Equation}
即
\begin{BoxEquation}[载流螺线管的磁场]
    B=\frac{\mu_0}{2}In(\cos\beta_2-\cos\beta_1)
\end{BoxEquation}
当螺线管的长度远远大于螺线管的半径，即$L\gg R$时，此时得到了一个无限长螺线管。

对于无限长螺旋管的内部各点处，有$\beta_1\rightarrow\pi$而$\beta_2\rightarrow 0$
\begin{BoxEquation}[载流无限长螺线管的轴线磁场]
    B=\mu_0nI
\end{BoxEquation}

\Refx{公式}{载流无限长螺线管的轴线磁场}指出，无限长载流螺线管轴线上的磁感强度与位置无关，也就是说，载流无限长螺线管在轴线上将产生匀强磁场，这是很重要的结论。\footnote{我们后面还会进一步证明，实际上无限长螺线管内部的磁场都是匀强磁场，不仅限于螺线管的轴线。}

对于半无限长螺线管的端点处，有$\beta_1=\pi/2$和$\beta_2\rightarrow 0$
\begin{BoxEquation}[载流半无限长螺线管的端点磁场]
    B=\frac{1}{2}\mu_0nI
\end{BoxEquation}
因而对于有限长的螺线管，其中央的磁感强度为$\mu_0nI$，其端点的磁感强度为$\mu_0nI/2$，即螺线管中央的磁感强度为螺线管端点处的磁感强度的两倍，且螺线管中间部分可以视为匀强磁场，当然以上这都是需要近似的结果，显然当$L$与$R$比值越大时近似程度越好。

我们先前在\Refx{公式}{载流螺线管的磁场}的推导中，由于三角换元的原因我们将结论转化为了基于$\beta_1,\beta_2$来表述，但很明显更为直观的办法通过场点的横坐标$x$表达。

我们为此将\Refx{图}{载流螺线管的磁场}的原点固定至螺线管中心。
\begin{Figure}{载流螺线管的中心坐标}
    \inputfigure[scale=0.8]{Chapter10B_Figure10}\vspace{-15pt}
\end{Figure}\vspace{-10pt}
由上图可以看出$x,\beta_1,\beta_2$的关系
\begin{Equation}
    \cos\beta_1=
    -\frac{x+L/2}{\sqrt{R^2+(x+L/2)^2}}\qquad
    \cos\beta_2=
    -\frac{x-L/2}{\sqrt{R^2+(x-L/2)^2}}
\end{Equation}
将上式代入\Refx{公式}{载流螺线管的磁场}
\begin{BoxEquation}[载流螺线管的磁场的坐标表示]
    B=\frac{\mu_0}{2}In\qty(\frac{x+L/2}{\sqrt{R^2+(x+L/2)^2}}-\frac{x-L/2}{\sqrt{R^2+(x-L/2)^2}})
\end{BoxEquation}
将\Refx{公式}{载流螺线管的磁场的坐标表示}绘制为函数图像
\begin{Figure}{载流螺线管的磁场分布图像}
    \inputfigure[scale=0.9]{Chapter10B_Figure09}
\end{Figure}
通过\Refx{图}{载流螺线管的磁场分布图像}我们可以直观的看到通电螺线管内的磁场分布，其中，红线是一个$L,R$比值更大的螺线管，蓝线是一个$L,R$比值更小的螺线管，由此我们可以看出$L,R$比值越大，中央区域越接近匀强磁场，两端磁感应强度的下降曲线越陡。另外也可以看到当$L,R$比值较小时，两端$\mu_0nI/2$的近似程度比中央$\mu_0nI$的近似程度好。

在本节首我们提到过，\textbf{载流螺线管的磁场与条形磁铁是等效的}，下图展现了这点
\begin{Figure}{载流螺线管的磁场与条形磁铁的对比}
    \subfigure[载流螺线管的磁场]
    {
        \includegraphics[scale=0.8]{old/Python/build/VFP_B04.pdf}
    }\qquad
    \subfigure[条形磁铁的磁场]
    {
        \includegraphics[scale=0.8]{old/Python/build/VFP_M04.pdf}
    }
\end{Figure}

\subsubsection{运动电荷的磁场}
在宏观上看，导体中的电荷是以体密度连续分布的，导体中的电流也是连续的，但是如果我们从微观上考察，电荷是若干离散的载流子，电流也就是若干载流子定向运动的结果。

在\Refx{定律}{毕奥-萨伐尔定律}中我们已经了解了电流产生的磁场，电流是载流子定向运动的结果，如果现在我们想通过电流的磁场推出运动电荷的磁场，那么，我们就有必要首先建立起\textbf{宏观连续的电流描述}和\textbf{微观离散的电荷运动描述}之间的关系，这是很有价值的。

第一组关系中，我们将讨论宏观和微观对电荷的不同观点。在宏观上我们认为电荷是连续分布的，可以通过单位体积的电荷量，即电荷体密度$\rho$来描述。在微观上电荷则是一系列离散的带电粒子，显然，带电粒子所带电荷量，带电粒子的分布越集中，其在宏观上就会表现处越大的电荷体密度，记单个电荷的带电量为$q$，记单位体积中的带电粒子数目为$n$，后者也称为\uwave{电荷数密度}或\uwave{电荷浓度}，那么我们就可以将宏观上某一处的电荷体密度$\rho$表示为
\begin{BoxEquation}[电荷的微观表述]
    \rho=nq
\end{BoxEquation}
第二组关系中，我们将讨论宏观和微观对电流的不同观点。为了简单起见，我们假设导体中的电荷和电流分布都是均匀的，若记导体垂直于电流的截面为$S$，那么顺次使用\Refx{公式}{电流与电流密度}和\Refx{定义}{电流密度}以及刚刚的\Refx{公式}{电荷的微观表述}
\begin{equation*}
    I=JS=\rho vS=nqvS
\end{equation*}

这一结果也可以从微观上解释，电流是单位时间内通过导体截面的电荷量，而能在单位时间内通过导体截面$S$的电荷必然分布在\textbf{导体截面前单位时间内电荷的运动距离内}，后者即带电粒子的速度，因而，这一区域的体积是$vS$，这一区域中的带电粒子数目即$nvS$，考虑到每个带电粒子的电荷量为$q$，因而这一区域内的总电荷量即$nqvS$，这也就是电流的大小。
\begin{BoxEquation}[电流的微观表述]
    I=nqvS
\end{BoxEquation}
第三步，我们要建立电流元和电荷运动的关系，我们知道电流元即$I\dif{\vb*{l}}$，其中$\dif{\vb*{l}}$的方向即电流的方向，它的方向等同于$q\vb*{v}$的方向，因此代入\Refx{公式}{电流的微观表述}时
\begin{equation*}
    I\dif{\vb*{l}}=nq\vb*{v}S\dif{l}
\end{equation*}

其中$S\dif{l}$即电流元包含的导体体积，可以记作$\dif{V}$
\begin{equation*}
    I\dif{\vb*{l}}=nq\vb*{v}\dif{V}
\end{equation*}

其中$n$为单位体积的带电粒子数目，故$n\dif{V}$即电流元中的带电粒子数目，记作$\dif{n}$
\begin{BoxEquation}[电流元与运动电荷]
    I\dif{\vb*{l}}=q\vb*{v}\dif{n}
\end{BoxEquation}
这样我们就完整在宏观和微观之间建立起了电流元与运动电荷的关系。现在让我们回到磁场的问题上，若以$\dif{\vb*{B}}$表示电流元的磁场，而以$\vb*{B}$表示电流元中每个运动电荷的磁场
\begin{Equation}[运动电荷的磁感强度1]
    \dif{\vb*{B}}=\vb*{B}nS\dif{l}=\vb*{B}n\dif{V}=\vb*{B}\dif{n}
\end{Equation}
这样就建立了电流元的磁场$\dif{\vb*{B}}$和运动电荷的磁场$\vb*{B}$的关系。

根据\Refx{定律}{毕奥-萨伐尔定律}和\Refx{公式}{电流元与运动电荷}
\begin{Equation}[运动电荷的磁感强度2]
    \dif{\vb*{B}}=
    \frac{\mu_0}{4\pi}
    \frac{I\dif{\vb*{l}}\times\vu*{r}}{r^2}=
    \frac{\mu_0}{4\pi}
    \frac{q\vb*{v}\times\vu*{r}}{r^2}\dif{n}
\end{Equation}
根据式\ref{式:运动电荷的磁感强度1}，即有
\begin{BoxEquation}[运动电荷的磁感强度]
    \vb*{B}=\frac{\mu_0}{4\pi}\frac{q\vb*{v}\times\vu*{r}}{r^2}
\end{BoxEquation}
\Refx{公式}{运动电荷的磁感强度}与\Refx{公式}{点电荷的电场强度}在形式上是对应的，而且也更为简单，那么为何我们最初要选取以电流元为中心的\Refx{定律}{毕奥-萨伐尔定律}作为磁场的起始呢？其中一个原因是，运动电荷构成的恒定电流可以形成静磁场，运动电荷的磁场却不是静磁场，因为作为磁场源的电荷的位置在不断变化，这违背了静磁学的要求。

但是，实际上一个更为重要的原因是，如果现在我们有两个电荷$q_1,q_2$分别以速度$\vb*{v}_1,\vb*{v}_2$在空间中运动，而$\vb*{r}$是$q_2$指向$q_1$的位矢，那么我们可以写出其中$q_1$受到的磁场力和电场力。

电荷$q_1$受到的磁场力$\vb*{F}_B$为
\begin{Equation}[电荷间的电场力]
    \vb*{F}_B=\frac{\mu_0}{4\pi}\frac{q_1q_2\vb*{v}_1\times\vb*{v}_2}{r^2}\times\vu*{r}
\end{Equation}
电荷$q_1$受到的电场力$\vb*{F}_E$为
\begin{Equation}[电荷间的磁场力]
    \vb*{F}_E=\frac{1}{4\pi\varepsilon_0}\frac{q_1q_2}{r^2}\vu*{r}
\end{Equation}
为了便于估计数量级，假定$\vb*{v}_1,\vb*{v}_2$垂直并且$\vb*{v}_2,\vb*{r}$平行，同时
\begin{equation*}
    q_1=q_2=1\si{C}\quad v_1=v_2=1\si{m\cdot s^{-1}}\quad r=1\si{m}
\end{equation*}
代入真空电容率$\varepsilon_0$和真空磁导率$\mu_0$的数值
\begin{equation*}
    \varepsilon_0=8.85\times 10^{-12}\si{C^2\cdot N^{-1}\cdot m^{-2}}\qquad
    \mu_0=4\pi\times 10^{-7}\si{N\cdot A^{-2}}
\end{equation*}
这就可以算得磁场力$F_B=4.0\times 10^{-7}$\si{N}，以及电场力$F_E=9.0\times 10^{9}$\si{N}。

这就是说，两个运动电荷间的电场力要比磁场力大了约$10^{16}$的数量级，因此，通过试探电荷的方式，我们绝无可能分析出另外一个运动电荷产生的磁场，我们会看到运动电荷间的相互作用与它们静止时毫无差别，因为电荷间的电场力要远远大于的磁场力，后者被遮蔽了。

所以\Refx{公式}{运动电荷的磁感强度}虽然很简洁，但其只是一个不能通过实验观测到的理论公式，而电流的磁场是可以实验观测的，因此其无法取代\Refx{定律}{毕奥-萨伐尔定律}作为磁场基本实验定律的地位，就像\Refx{定律}{库仑定律}作为电场基本实验定律那样。

所以为何运动电荷间的磁场力比电场力弱如此之多？以及为何运动电荷间的磁场力无法被观测，但是大量运动电荷构成的电流的磁场就可以被观测？下一小节将解答这两个问题。

\subsection{磁场与电场的统一性}
这一小节让我们来思考一个有趣的问题，我们在式\ref{式:磁场的作用}中曾提及运动电荷受到的磁场力可以表示为$\vb*{F}=q\vb*{v}\times\vb*{B}$，我们可能会疑惑这里的电荷速度$\vb*{v}$是相对于哪一参照系的？

关于选取什么样的参照系作为规定磁场的合适参照系，我们还未曾说过什么。而事实是，任何一个惯性系都可以！对此我们立即会表达出以下的质疑，电荷在不同的惯性参照系中将具有不同的速度，那么按照这里的观点，电荷将受到不同的磁场力，这也就是说电荷受到的磁场力的多少竟然会取决于我们选取的参照系的选取，这听上去实在是有些荒唐。

下面这个例子可以比较清楚的反映这个问题，设有一根电流为$I$在载流长直导线，设导线内的电子以速度$\vb*{v}$向右运动，设导线外有一带负电的粒子同样以速度$\vb*{v}$向右运动。\vspace{5pt}
\begin{Figure}{磁场与电场的统一性}
    \subfigure[基于导线的参照系]
    {
        \inputfigure[scale=0.85]{Chapter10B_Figure11}
    }\\ \vspace{10pt}
    \subfigure[基于电荷的参照系]
    {
        \inputfigure[scale=0.85]{Chapter10B_Figure12}
    }
\end{Figure}
如果我们以导线为参照系，根据$\vb*{F}=qv\times\vb*{B}$可知，此时导线外的电荷将会受到一个指向导线的磁场力，但是，如果我们以电荷为参照系，此时电荷自身的速度变为了零，而静止电荷是不会受到任何磁场力的，然而物体的受力又不应随着参照系的选取而变，因此电荷此时必然将受到一个指向导线的电场力，但是这就意味着运动的“中性”载流导线将带正电。

我们有必要对此仔细检查，为何载流导线在运动时将带电？首先需要肯定的是，电荷量不会随着运动而变化。电荷量的微小差异就会产生巨大的电场力，我们知道物质是中包含带正电的质子和带负电的电子，当我们将物质加热时，由于质子和电子的质量不同，两者因热运动取得的速度增量是不同的，倘若电荷量真与运动有关，两者的电荷量就无法平衡了，而且这种不平衡将会在宏观上可观测，但是我们却从未观察到某种物质在加热后会带电。

我们尚还没有系统的掌握相对论，但解释这里的现象，我们只需要明白以下这样一个简单的事实，\textbf{运动的东西会在运动方向上变得更短}，即\uwave{相对论收缩}，或更为熟知的\uwave{尺缩效应}。

我们知道，在金属导体中，带有负电荷的电子通过运动传递电流，带有正电荷的金属阳离子在材料中是保持不动的，而静止的导线显然是电中性的，即在导线参照系中有$\rho_{+}=\rho_{-}$。

我们现在来考察正负电荷在参照系变换时究竟发生了什么
\begin{itemize}
    \item 以导线为参照系时，正电荷是静止的，以运动电荷为参照系时，正电荷向左运动，根据尺缩原理，新的参照系中正电荷的间距会变小，正电荷的密度$\rho_{+}$会变得更大。
    \item 以导线为参照系时，负电荷向右运动，以运动电荷为参照系时，负电荷是静止的，根据尺缩原理，新的参照系中正电荷的间距会变大，正电荷的密度$\rho_{+}$会变得更小。
\end{itemize}
这就说明在运动电荷的参照系下，这根运动的导线中$\rho_{+}>\rho_{-}$，因而将带有正净电荷，从而对导线外的电荷产生了一个指向导线的电场力，这个力的大小等同于原先的磁场力。

这个例子就告诉我们，电场和磁场是不可分裂或者说是不可分辨的，也就是说，电场和磁场是一个统一整体的两个不同侧面，这个统一整体也就是我们所谓的\uwave{电磁场}。电和磁之间没有严格的界限，电磁场在具体问题中究竟以何种比例表现为电场和磁场，仅取决于我们选取的参照系，就像在这个例子中，在导线参照系中就仅有磁场，在电荷参照系中就仅有电场。

电磁场在我们熟悉的参照系中主要表现为电场，从这样一个观点看，磁场可以被认为是电场由于电荷的运动产生的相对论效应，这种相对论修正的数量级往往在$v^2/c^2$左右，这也是为何先前我们看到运动电荷间的电场力是磁场力的$10^{16}$倍左右，因为光速$c=3\times 10^8$\si{m\cdot s^{-1}}。

磁场是电场的相对论效应，但是，且慢！我们知道相对论效应实际上只有在高速时才会较为显著，在低速时是完全可以忽略的，那么磁场作为一种相对论效应，理论上在我们的低速世界中是不应该被观测到的。然而，实际上我们却可以观测到许多磁效应，例如两根平行载流导线间的磁相互作用力是可以通过实验观测的。磁场作为一种$v^2/c^2$级别的相对论修正对于这里导线中低于$1$\si{m\cdot s^{-1}}的电荷来说本应是可以忽略的，但电流的特殊之处在于，在导线中正负电荷被极为精确的抵消，以至于正常的电场完全消失了，这才使得磁场可以被观测。
\begin{center}
    \small 
    正是电流的这种特殊性质，我们才能在低速世界观测到丰富的磁现象和磁效应。
\end{center}
这也就解释了为何电流的磁场可以被实验观测，而运动电荷的磁场则不行。

\vspace{-10pt}
\section{磁场高斯定理与磁通量}
\vspace{-10pt}

\subsection{磁感线}
为了形象的描述空间磁场的分布，我们可以引入\uwave{磁感线}的概念，简称$\vb*{B}$线。

磁感线是这样的一簇曲线，每条磁感线上任意一点的切线方向，即代表该点处的电场方向，同时，磁感线越密集则代表磁感强度越大，磁感线线越稀疏代表电场强度越小，简单的说
\begin{center}
    磁感线的方向反映电场的方向，磁感线的疏密反映电场的强弱。
\end{center}
磁感线和电场线均是用于描绘场分布的场线，既有相似之处，又有不同之处。

磁感线和电场线的共同特点是两者的场线都不会出现相交的情形，这是因为场线的方向反映矢量场的方向，如果场线相交就意味着矢量函数在交点处有两个方向，这是不可能的。

磁感线与电场线的不同之处在于以下的特征
\begin{itemize}
    \item 磁感线总是形成闭合曲线，不存在起始点，不存在终止点。
    \item 电场线不会形成闭合曲线，起始于正电荷，终止于负电荷。
\end{itemize}
我们知道，场线是否有起始点和终止点的特性反映了场是否有源，场线是否形成闭合曲线反映了场是否有旋。因此磁感线和电场线特征上的差异，就表明磁场和电场实际上是两种性质不同的场，磁场是有旋无源场，电场是有源无旋场，有关电场源和旋的性质我们已经在先前的第八章中证明了，有关磁场源和旋的性质将在本节和下一节分别予以证明。

下表将列举几种常见电流产生的磁场的磁感线分布，为了更好的比较磁感线与电场线异同，以及其表征的磁场和电场的异同，下表将同时列出与之相似的电场的电场线分布。\footnote{实际上，不同方向的载流导线应该对应带不同电荷的带电细棒，而不是点电荷，当然这不是很要紧。}\vspace{20pt}
\begin{TableLong}{常见磁感线与电场线的对比}
{|>{\small}c|>{\small}c|}
{}
    \includegraphics[scale=0.8]{old/Python/build/VFP_B02.pdf}
    &
    \includegraphics[scale=0.8]{old/Python/build/VFP_E01.pdf}\\ \hline
    垂直纸面向外的电流&正点电荷\\ \hline
    \includegraphics[scale=0.8]{old/Python/build/VFP_B01.pdf}
    &
    \includegraphics[scale=0.8]{old/Python/build/VFP_E02.pdf}\\ \hline
    垂直纸面向内的电流&负点电荷\\ \hline
    \includegraphics[scale=0.8]{old/Python/build/VFP_B03.pdf}
    &
    \includegraphics[scale=0.8]{old/Python/build/VFP_E03.pdf}\\ \hline
    磁偶极子&电偶极子\\ \hline
    \includegraphics[scale=0.8]{old/Python/build/VFP_B04.pdf}
    &
    \includegraphics[scale=0.8]{old/Python/build/VFP_E04.pdf}\\ \hline
    通电螺线管&带电平面组\\ \hline
\end{TableLong}

载流直导线和载流圆线圈轴线处的磁场都可以通过右手螺旋法则简记
\begin{itemize}
    \item 如果以右手伸直的拇指表示电流方向，即载流直导线，那么弯曲的四指表示磁场方向。
    \item 如果以右手弯曲的四指表示电流方向，即载流圆线圈，那么伸直的拇指表示磁场方向。
\end{itemize}
% 载流螺线管的情况与载流圆线圈类似的，其管内的磁场也可以通过右手螺旋法则判定。

在上表中我们看到，磁偶极子和电偶极子产生的场虽然在靠近场源电流和场源电荷的部分表现的非常不同，但两者在远离场源的远场点处则是非常接近的，这就验证了先前的说法。

在上表中我们也看到，磁场中的\textbf{无限长通电螺线管}和电场中的\textbf{无限大带电平面组}具有某种相似性，前者可以产生匀强磁场，后者可以产生匀强电场，两者的地位是相当的，我们在稍后的内容中还将看到电路中与\textbf{电容}对应的另一种重要元器件\textbf{电感}也就是由螺线管构成的。\vspace{-10pt}

\subsection{磁通量}
定义磁感强度$\vb*{B}$在某一给定曲面$S$上的通量为\uwave{磁通量}，用$\Phi_B$表示。
\begin{BoxDefinition}[磁通量]
    \Phi_B=\Isnt[S]{\vb*{B}\cdot\dif{\vb*{S}}}
\end{BoxDefinition}
磁通量在国际单位制中的单位是\si{Wb=T\cdot m^2}，称为\uwave{韦伯}，量纲是\si{[L^2MT^{-2}I^{-1}]}。

磁通量$\Phi_B$与先前的电通量$\Phi_E$是对应的，两者的性质和意义是基本相似的。\vspace{-10pt}

\subsection{磁场高斯定理}
在先前的\Refx{定律}{电场高斯定理}中我们研究了闭合曲面上的电通量，在这里我们就要对应的研究闭合曲面$S$上的磁通量，假定空间中仅有一根电流为$I$的载流导线$l'$。

根据\Refx{定律}{毕奥-萨伐尔定律}，载流导线$l'$产生的磁场为
\begin{Equation}[磁场高斯定理0]
    \vb*{B}=\Ilnt[l']{\frac{\mu_0}{4\pi}\frac{I\dif{\vb*{l'}\times\vu*{r}}}{r^2}}
\end{Equation}
根据高斯公式，若欲求$S$上的磁通量，就要先求出$\vb*{B}$的散度。
\begin{Equation}[磁场高斯定理1]
    \Isot[S]{\vb*{B}\cdot\dif{\vb*{S}}}=\Itnt{\nabla\cdot\vb*{B}\dif{V}}
\end{Equation}
式\ref{式:磁场高斯定理1}中的$\Omega$是闭曲面$S$所包围的空间区域，代入式\ref{式:磁场高斯定理0}得
\begin{Equation}[磁场高斯定理2]
    \Isot[S]{\vb*{B}\cdot\dif{\vb*{S}}}=
    \Ilnt[l']{\frac{\mu_0}{4\pi}
    \Itnt{\nabla\cdot\qty(\frac{I\dif{\vb*{l'}\times\vu*{r}}}{r^2})\dif{V}}}
\end{Equation}
根据向量微分算子的运算法则
\begin{Equation}[磁场高斯定理3]
    \nabla\cdot(\vb*{A}\times\vb*{B})=
    (\nabla\times\vb*{A})\cdot\vb*{B}-
    (\nabla\times\vb*{B})\cdot\vb*{A}
\end{Equation}
这里就有
\begin{Equation}[磁场高斯定理4]
    \nabla\cdot\qty(\frac{I\dif{\vb*{l'}\times\vu*{r}}}{r^2})=
    \qty(\nabla\times I\dif{\vb*{l'}})\cdot\frac{\vu*{r}}{r^2}-
    \qty(\nabla\times\frac{\vu*{r}}{r^2})\cdot I\dif{\vb*{l'}}
\end{Equation}
式\ref{式:磁场高斯定理4}的第一项中的$\nabla\times I\dif{\vb*{l'}}$所涉及的$I\dif{\vb*{l'}}$对于曲线积分中每一取定的电流元来说是一个常矢量，而常矢量的旋度处处为零矢量，因此第一项为零，我们不必再考虑了。

式\ref{式:磁场高斯定理4}的第二项可以这样计算，关注到
\begin{equation*}
    \frac{\vu*{r}}{r^2}=\frac{\vb*{r}}{r^3}=\frac{x\vi+y\vj+z\vk}{(x^2+y^2+z^2)^{\frac{3}{2}}}=P\vi+Q\vj+R\vk
\end{equation*}
考虑到\vspace{-10pt}
\begin{align*}
    \Fpif{R}{y}=\frac{\frac{3}{2}z(x^2+y^2+z^2)^{\frac{1}{2}}2y}{(x^2+y^z+z^2)^{3}}=\frac{3zy}{(x^2+y^2+z^2)^{\frac{5}{2}}}\\[3mm]
    \Fpif{Q}{z}=\frac{\frac{3}{2}y(x^2+y^2+z^2)^{\frac{1}{2}}2z}{(x^2+y^z+z^2)^{3}}=\frac{3yz}{(x^2+y^2+z^2)^{\frac{5}{2}}}
\end{align*}
这就可以得出
\begin{equation*}
    \Fpif{R}{y}-\Fpif{Q}{z}=0
\end{equation*}
类似的还有
\begin{equation*}
    \Fpif{P}{z}-\Fpif{R}{x}=0\qquad
    \Fpif{Q}{x}-\Fpif{P}{y}=0
\end{equation*}
因此就有$\nabla\times\dfrac{\vu*{r}}{r^2}=\vb*{0}$，从而式\ref{式:磁场高斯定理4}的第二项也为零。

那么式\ref{式:磁场高斯定理4}的结果就为
\begin{equation*}
    \nabla\cdot\qty(\frac{I\dif{\vb*{l'}\times\vu*{r}}}{r^2})=0
\end{equation*}

那么将上式代入式\ref{式:磁场高斯定理2}，就可以得到
\begin{equation*}
    \Isot[S]{\vb*{B}\cdot\dif{\vb*{S}}}=0
\end{equation*}
以上的推导就说明，载流导线对闭合曲面的磁通量贡献总是零，因此
\begin{BoxPrinciple}[磁场高斯定理]
    在磁场中，通过任意一个闭合曲面的磁通量总是为零。
    \begin{BoxEq}
        \Isot[S]{\vb*{B}\cdot\dif{\vb*{S}}}=0
    \end{BoxEq}
\end{BoxPrinciple}
该规律就称为\uwave{磁场高斯定理}，在定理中所取的闭合曲面$S$称为\uwave{高斯面}。

\Refx{定律}{磁场高斯定理}也可以通过磁感线形象理解，
如果磁感线完全位于闭曲面内部或外部，那其对闭曲面的磁通量就没有任何贡献，如果磁感线穿过了闭曲面，由于磁感线不能中断，其从一处穿入闭曲面后必然会从另一处穿出，两者的磁通量相互抵消。

\Refx{定律}{磁场高斯定理}也可以使用微分形式表达为
\begin{BoxPrinciple}[磁场高斯定理的微分形式]
    在磁场中，磁感强度的散度恒为零。
    \begin{BoxEq}
        \nabla\cdot\vb*{B}=0
    \end{BoxEq}
\end{BoxPrinciple}
\Refx{定律}{磁场高斯定理的微分形式}指出
\begin{center}
    磁场是无源场
\end{center}
磁场是无源场，这就可以解释描绘磁场分布的磁感线为何没有起始点和终止点了。\vspace{10pt}

\section{磁场的安培环路定理}
静电场是一个保守场，因此\Refx{定律}{电场环路定理}指出静电场$\vb*{E}$的环流为零，那么现在的问题就是，静磁场是否一个保守场？事实上并不是，静磁场是非保守的，静磁场中也因此不能定义所谓的磁势的概念，当然,这也就意味着静磁场具有不为零的$\vb*{B}$的环流。

本节我们就来研究静磁场的环流与什么因素有关，就像先前研究静电场的通量一样。

\subsection{安培环路定理}
在空间中任意取一条有向闭合曲线$L$，假定空间中仅有一根电流为$I$的载流导线$l'$。

根据\Refx{定律}{毕奥-萨伐尔定律}，载流导线$l'$产生的磁场为
\begin{Equation}[安培环路定律0]
    \vb*{B}=\Ilnt[l']{\frac{\mu_0}{4\pi}\frac{I\dif{\vb*{l'}\times\vu*{r}}}{r^2}}
\end{Equation}
根据斯托克斯公式，若欲求$L$上的磁场环流，就要先求出$\vb*{B}$的旋度。
\begin{Equation}[安培环路定律1]
    \Ilot[L]{\vb*{B}\cdot\dif{\vb*{l}}}=\Isnt[\Sigma]{\nabla\times\vb*{B}\cdot\dif{\vb*{S}}}
\end{Equation}
式\ref{式:安培环路定律1}中的$\Sigma$是闭曲线$L$所张成的任意空间曲面，代入式\ref{式:安培环路定律0}得
\begin{Equation}[安培环路定律2]
    \Ilot[L]{\vb*{B}\cdot\dif{\vb*{l}}}=
    \Ilnt[l']{\frac{\mu_0}{4\pi}
    \Isnt{\nabla\times\qty(\frac{I\dif{\vb*{l'}\times\vu*{r}}}{r^2})\cdot\dif{\vb*{S}}}}
\end{Equation}
根据向量微分算子的运算法则\footnote{这里的$(\vb*{A}\cdot\nabla)=A_x\Fpif{x}+A_y\Fpif{y}+A_z\Fpif{z},~(\vb*{B}\cdot\nabla)=B_x\Fpif{x}+B_y\Fpif{y}+B_z\Fpif{z}$，可以视为一个完整的微分算子。}
\begin{Equation}[安培环路定律3]
    \qquad\qquad\qquad\qquad
    \nabla\times(\vb*{A}\times\vb*{B})=
    \vb*{A}(\nabla\cdot\vb*{B})-
    \vb*{B}(\nabla\cdot\vb*{A})+
    (\vb*{B}\cdot\nabla)\vb*{A}-
    (\vb*{A}\cdot\nabla)\vb*{B}
    \qquad\qquad\qquad\qquad
\end{Equation}
这里就有
\begin{Equation}[安培环路定律4]
    \qquad
    \nabla\times\qty(\frac{I\dif{\vb*{l'}\times\vu*{r}}}{r^2})=
    I\dif{\vb*{l'}}\qty(\nabla\cdot\frac{\vu*{r}}{r^2})-
    \frac{\vu*{r}}{r^2}\qty(\nabla\cdot I\dif{\vb*{l'}})+
    \qty(\frac{\vu*{r}}{r^2}\cdot\nabla)I\dif{\vb*{l'}}-
    \qty(I\dif{\vb*{l'}}\cdot\nabla)\frac{\vu*{r}}{r^2}\qquad
\end{Equation}
式\ref{式:磁场高斯定理4}第二项和第三项均是对常矢量$I\dif{\vb*{l'}}$的微分运算，其结果都为零。

式\ref{式:磁场高斯定理4}第四项并不是零，但是有一些文献认为，如果我们考虑的载流导线$l'$本身也是一个回路，这一项将在环路积分中通过全微分求积的方式被消掉，然而这中间的数学原理尚未被弄清，因此这里我们先不加证明的承认第四项的积分结果为零，以继续我们的推导。

那么现在我们只需要处理第一项
\begin{equation*}
    \frac{\vu*{r}}{r^2}=\frac{\vb*{r}}{r^3}=\frac{x\vi+y\vj+z\vk}{(x^2+y^2+z^2)^{\frac{3}{2}}}=P\vi+Q\vj+R\vk
\end{equation*}
考虑到\vspace{-10pt}
\begin{align*}
    \Fpif{P}{x}=\frac{(x^2+y^2+z^2)-3x^2}{(x^2+y^2+z^2)^{\frac{5}{2}}}\\[3mm]
    \Fpif{Q}{y}=\frac{(x^2+y^2+z^2)-3y^2}{(x^2+y^2+z^2)^{\frac{5}{2}}}\\[3mm]
    \Fpif{R}{z}=\frac{(x^2+y^2+z^2)-3z^2}{(x^2+y^2+z^2)^{\frac{5}{2}}}
\end{align*}
这就可以得出
\begin{equation*}
    \Fpif{P}{x}+\Fpif{Q}{y}+\Fpif{R}{z}=0
\end{equation*}
因此就有$\nabla\cdot\dfrac{\vu*{r}}{r^2}=\vb*{0}$，即
\begin{equation*}
    \nabla\times\qty(\frac{I\dif{\vb*{l'}\times\vu*{r}}}{r^2})=0
\end{equation*}
若载流导线没有穿过闭合曲线，那么总有$r\neq 0$，将上式代入\ref{式:安培环路定律2}
\begin{Equation}[安培环路定律6]
    \Ilot[L]{B\cdot\dif\vb*{l}}=0
\end{Equation}
若载流导线穿过闭合曲线，那么此时无论以闭合曲线为边界的空间曲面$\Sigma$如何选取，这里的载流导线$l'$总是存在一个与空间曲面$\Sigma$的交点，而在交点上有$r=0$，这会使得上方三组偏导数的分母为零，构成奇点，导致偏导数不连续，此时斯托克斯公式无法使用。

为了继续使用斯托克斯公式，我们在载流导线附近取一个半径很小且垂直于载流导线的有向圆环$L_0$，那么此时闭曲线$L$和圆环$L_0$之间的空间曲面$\Sigma_R$并不包含$r=0$的点。

因此可以在$\Sigma_R$上应用斯托克斯公式，其边界曲线为$L$和$L_0$，从而式\ref{式:电场高斯定理1}被修正为\footnote{这里设$L_0$的绕向与$L$基本一致，但是$L_0$作为$\Sigma_R$的内边界曲线应取相反的绕向，故其曲线积分前有负号。}
\begin{Equation}[安培环路定律5]
    \Ilot[L]{B\cdot\dif\vb*{l}}-\Ilot[L_0]{B\cdot\dif\vb*{l}}=0
\end{Equation}
那么现在的问题，就是要求出圆环$L_0$上的磁场环流。

这里由于圆环$L_0$取的很小，这时的载流导线$l'$相对于$L_0$可以视为一条垂直通过圆环$L_0$中心的无限长直导线，设圆环$L_0$的半径为$r$，根据\Refx{公式}{载流无限长直导线的磁场}
\begin{equation*}
    \Ilot[L_0]{B\cdot\dif\vb*{l}}=
    \Ilot[L_0]{\frac{\mu_0}{2\pi}\frac{I}{r}\dif{l}}=\frac{\mu_0}{2\pi}\frac{I}{r}2\pi r=\mu_0I
\end{equation*}

这里圆环$L_0$和闭曲线$L$的绕向是一致的，如果假定电流$I$的正值表示$L$的右手螺旋方向
\begin{itemize}
    \item 当电流$I$为正时，磁场方向与$L_0$的绕向相同，磁场环流为正，而$\mu_0I$恰好也为正。
    \item 当电流$I$为负时，磁场方向与$L_0$的绕向相反，磁场环流为负，而$\mu_0I$恰好也为负。
\end{itemize}
因此上式总是正确的，将其代入\ref{式:安培环路定律6}即得
\begin{BoxPrinciple}[安培环路定理]
    在磁场中，沿着任意一个闭合曲线的磁场环流，

    只与闭曲线所包围的电流的大小有关，而与闭曲线外的电流分布无关。
    \begin{BoxEq}
        \Ilot[L]{\vb*{B}\cdot\dif{\vb*{l}}}=\mu_0I
    \end{BoxEq}
    其中，通过闭曲线的电流以闭合曲线的右手螺旋方向为正值。
\end{BoxPrinciple}
该规律就称为\uwave{安培环路定理}，在定理中所取的闭合曲线$L$称为\uwave{安培环路}。

根据旋度的定义，磁场在某一点处的旋度，是该点处的磁场环流密度，即\footnote{其中$P$为磁场中某一点，而$\Sigma$表示闭曲线$L$张成的空间曲面，而$S$表示空间曲面$\Sigma$的面积。}
\begin{equation*}
    \nabla\times\vb*{B}=\Lim{\Sigma}{P}{\frac{1}{S}\Ilot[L]{\vb*{B}\cdot\dif{\vb*{l}}}}
\end{equation*}
对于线状电流形成的磁场而言，当空间曲面$\Omega$连同其边界曲线$L$缩向点$P$时，若点$P$不在导线上，那么该点处的旋度为零，若点$P$在导线上，那么围绕该点的闭曲线上的磁场环流为定值$\mu_0I$，当该闭曲面缩向一点时，随着$S\rightarrow 0$，可得\textbf{线状电流处的旋度为无穷大}。

因此线状电流形成的磁场是无法讨论磁场的旋度的，线状电流的模型不适用该情境。

但是如果我们考虑电流连续分布的情况，例如考虑到导线的截面宽度，由于
\begin{equation*}
    \nabla\times\vb*{B}=\Lim{\Sigma}{P}{\frac{1}{S}\Ilot[L]{\vb*{B}\cdot\dif{\vb*{l}}}}=\Lim{\Sigma}{P}{\frac{\mu_0I}{S}}=\mu_0\vb*{J}
\end{equation*}
这就得到了一个重要定律
\begin{BoxPrinciple}[安培环路定理的微分形式]
    在磁场中，磁感强度的旋度，与该点处的电流密度成正比。
    \begin{BoxEq}
        \nabla\times\vb*{B}=\mu_0\vb*{J}
    \end{BoxEq}
\end{BoxPrinciple}
\Refx{定律}{安培环路定理的微分形式}指出
\begin{center}
    磁场是有旋场
\end{center}
磁场是有旋场，这就可以解释为何磁感线总是表现为涡旋的闭合曲线。

\Refx{定律}{安培环路定理的微分形式}指出，如果空间中某处没有电流分布，由于此处电流密度为零，那么该处磁感强度的旋度也为零，即磁场在没有电流分布处无旋，或者说
\begin{center}
    磁场的旋完全来自于电流的存在
\end{center}
至此，我们可以归纳出以下的磁场基本性质
\begin{BoxPrinciple}[磁场的基本性质]
    磁场是一个有旋无源的非保守场。
\end{BoxPrinciple}

\subsection{安培环路定理的应用}
现在我们来通过安培环路定理计算某些具有一定的对称性的载流导线的磁场分布，就像先前我们用电场高斯定理可以计算某些具有一定的对称性的带电体的电场分布那样。

\subsubsection{无限长均匀圆柱载流导体的磁场}
现在考虑一根半径为$R$的长直圆柱导体，其截面$S=\pi R^2$上均匀的载有电流$I$，根据对称性的原理，这跟圆柱导体上的电流产生的磁场必然是以圆柱轴线为中心的柱对称分布，这也就是说，如果我们记距离轴线为$r$处的磁感强度为$B$，并且记$r$处的圆周为$L$，那么
\begin{equation*}
    \Ilot[L]{\vb*{B}\cdot\dif{\vb*{l}}}=2\pi rB
\end{equation*}
如果$r<R$，即安培环路$L$取在导体内，那么此时$L$内包含的电流$I_1$
\begin{equation*}
    I_1=\frac{\pi r^2}{\pi R^2}I=\frac{r^2}{R^2}I
\end{equation*}
如果$r>R$，即安培环路$L$取在导体内，那么此时$L$内包含的电流即为$I_2=I$。

根据\Refx{定律}{安培环路定理}，当$r<R$时
\begin{equation*}
    \Ilot[L]{\vb*{B}\cdot\dif{\vb*{l}}}=\mu_0I_1=\mu_0\frac{r^2I}{R^2}
\end{equation*}
因而就有
\begin{Equation}[无限长圆柱载流导体的磁场1]
    B=\frac{\mu_0}{2\pi}\frac{Ir}{R^2}\qquad (r<R)
\end{Equation}
根据\Refx{定律}{安培环路定理}，当$r>R$时
\begin{equation*}
    \Ilot[L]{\vb*{B}\cdot\dif{\vb*{l}}}=\mu_0I_1=\mu_0I
\end{equation*}
因而就有
\begin{Equation}[无限长圆柱载流导体的磁场2]
    B=\frac{\mu_0}{2\pi}\frac{I}{r}\qquad (r>R)
\end{Equation}
联立式\ref{式:无限长圆柱载流导体的磁场1}和式\ref{式:无限长圆柱载流导体的磁场2}，即有
\begin{BoxEquation}[均匀圆柱载流导体的磁场]
    B=
    \begin{cases}
        \dfrac{\mu_0}{2\pi}\dfrac{Ir}{R^2},&r<R\\[5mm]
        \dfrac{\mu_0}{2\pi}\dfrac{I}{r},&r>R
    \end{cases}
\end{BoxEquation}
这一结论指出，均匀圆柱载流导体内部的磁场与半径成正比，均匀圆柱载流导体外部的磁场与载流细导线的磁场是一致的，与半径成反比，这也是为何通常可以忽略导线的截面面积。

\subsubsection{通电螺线管的磁场}
先前通过\Refx{定律}{毕奥-萨伐尔定律}，我们已经在\Refx{公式}{载流无限长螺线管的轴线磁场}中看到无限长载流螺线管轴线上的磁场是匀强的，现在我们将进一步通过\Refx{定律}{安培环路定理}证明，螺线管内部的磁场均为匀强的，螺线管外部则没有磁场。
\begin{Figure}{证明螺旋管内外的磁场方向}
    \subfigure[原始螺线管]
    {
        \label{图:原始螺线管}
        \inputfigure[scale=0.45]{Chapter10C_Figure01}
    }
    \subfigure[将螺线管翻转]
    {
        \label{图:将螺线管翻转}
        \inputfigure[scale=0.45]{Chapter10C_Figure02}
    }
    \subfigure[将电流反转]
    {
        \label{图:将电流反转}
        \inputfigure[scale=0.45]{Chapter10C_Figure03}
    }\vspace{-10pt}
\end{Figure}
首先我们来证明，螺线管内任意一点的磁场方向都为平行于轴线的方向，假定原先螺线管内场点$P$处的磁感强度有两种可能的情况，可能平行于轴线$\vb*{B}$，可能不平行于轴线$\vb*{B}'$。

第一种变换，由于螺线管无限长，因此任何一个场点$P$都可以视为位于螺线管轴线的中心位置，因而我们可以将其沿过点$P$的对称轴$y$轴翻转，此时磁感强度将会左右翻转。

第二种变换，将螺线管上的电流反转方向，此时磁感强度将会变为原先的反方向。

然而我们不难看出，这两种变换的结果是一致的，那么磁感强度的结果理应也是一致的，即应有$\vb*{B}_1=\vb*{B}_2$和$\vb*{B}_1'=\vb*{B}_2'$，但是我们从图中可以看出，后者是不成立的，因此螺线管内的磁场方向必须平行于轴线，这一反证过程也适用于螺线管外的场点，也就是说，如果螺线管外存在磁场，那么其也是平行于螺线管的轴线的（后面我们会看到螺线管外实际没有磁场）。

另外我们还可以得出一个结论，任何一条平行于螺线管轴线的直线上的各点，其不仅具有相同的磁场方向，其磁场的大小也相等，因为我们并不能区分无限长螺线管在左右上的位置。

随后我们来证明，螺线管内的磁场与轴线相同，螺线管外的磁场则总是为零。
\begin{Figure}{证明螺线管内为匀强磁场}
    \inputfigure[scale=0.95]{Chapter10C_Figure04}
\end{Figure}
在螺线管内取一个矩形环路$L_a$，其中$L_a$的$a_1$段在螺线管的轴线上，其方向与磁场方向保持一致，其长度为$l$，关注到$a_2,a_3$段与磁场垂直，因此磁场在$a_2,a_3$段上的积分为零。

在矩形环路$L_a$内没有包含任何电流，因此根据\Refx{定律}{安培环路定理}
\begin{equation*}
    \Ilot[L_1]{\vb*{B}\cdot\dif{\vb*{l}}}=Bl-B_1l=0
\end{equation*}
在上式中可以解得
\begin{equation*}
    B_1=B=\mu_0nI
\end{equation*}
这就说明螺线管内任意一处磁场均与轴线上的磁场相同，即螺线管内为匀强磁场。类似的我们也可以证明螺线管外的情况，只不过此时矩形环路$L_b$将包含电流$\mu_0nlI$，故
\begin{equation*}
    \Ilot[L_2]{\vb*{B}\cdot\dif{\vb*{l}}}=Bl-B_2l=\mu_0nlI
\end{equation*}
从而就有
\begin{equation*}
    B_2=0
\end{equation*}
这就说明螺线管外的磁场总是为零。
\begin{BoxPrinciple}[螺线管的磁场分布]
    对于一个无限长载流螺线管，其内部为线状匀强磁场，其外部无磁场。
    \begin{BoxEq}
        B=
        \begin{cases}
            \mu_0nI,&\text{\footnotesize 螺线管内部}\\
            0,&\text{\footnotesize 螺线管外部}
        \end{cases}    
    \end{BoxEq}
\end{BoxPrinciple}

\subsubsection{通电螺绕环}
这里我们来考虑一个新的结构，我们知道，将导线密绕在柱面上就会得到一个螺线管，非常类似的，将导线密绕在环面上就会得到一个\uwave{螺绕环}，它的剖面图如下图所示。
\begin{Figure}{螺绕环与螺线管}
    \subfigure[螺绕环]
    {
        \inputfigure[scale=0.8]{Chapter10C_Figure05}
    }\qquad\qquad
    \subfigure[螺线管]
    {
        \inputfigure[scale=0.8]{Chapter10C_Figure06}
    }
\end{Figure}

螺绕环在某种意义上可以视为一个“掰弯”后“首尾相接”的螺线管，因此可以预期螺绕环的性质与螺线管是相似的。设螺绕环的内外径为$R_1,R_2$，现在考虑螺绕环内部半径为$r$处的磁场，根据对称性的原理，这里的磁场是柱对称的，因而对于取在$r$处的安培环路$L$
\begin{equation*}
    \Ilot[L]{\vb*{B}\cdot\dif{\vb*{l}}}=2\pi rB
\end{equation*}
螺绕环的线圈密度若记为$n$，那么安培环路$L$就包含了$2\pi R_1nI$的电流。

根据\Refx{定律}{安培环路定理}
\begin{equation*}
    \Ilot[L]{\vb*{B}\cdot\dif{\vb*{l}}}=2\pi\mu_0 R_1nI
\end{equation*}
这就可以解得
\begin{equation*}
    B=\mu_0\frac{R_1}{r}nI
\end{equation*}
这里为了简便起见，我们假定螺绕环中的圆环是很细的，那么便有$r=R_1=R_2$，从而
\begin{equation*}
    B=\mu_0nI
\end{equation*}
这样就证明了螺绕环内部为与螺线管相同的匀强磁场。通过\Refx{定律}{安培环路定理}也可以类似证明螺绕环外的磁场为零，取在$R_1$之内的安培环路中没有电流，取在$R_2$之外的安培环路中，内径上的电流和外径上的电流相互抵消了，因此计算得出的磁场均为零。
\begin{BoxPrinciple}[螺绕环的磁场分布]
    对于一个很细的载流螺绕环，其内部为环状匀强磁场，其外部无磁场。
    \begin{BoxEq}
        B=
        \begin{cases}
            \mu_0nI,&\text{\footnotesize 螺绕环内部}\\
            0,&\text{\footnotesize 螺绕环外部}
        \end{cases}    
    \end{BoxEq}
\end{BoxPrinciple}


% 在导体外取一半径为$r$的安培环路$L_1$，根据\Refx{定律}{安培环路定律}
% \begin{equation*}
%     \Ilot[L]{\vb*{B}\cdot\dif{\vb*{l}}}=\mu_0I
% \end{equation*}

\section{带电粒子在磁场中的运动}

\subsection{洛伦兹力}
在本章初我们曾经提到过，阴极射线管中的电子束在磁场力的作用下将发生偏转，本节将讨论单个点电荷在运动时所受的磁场力，这种运动电荷在磁场中受到的力称为\uwave{洛伦兹力}
\begin{BoxEquation}[洛伦兹力]
    \vb*{F}=q\vb*{v}\times\vb*{B}
\end{BoxEquation}
\Refx{公式}{洛伦兹力}是实验结论，其表明，运动电荷受到的磁场力的方向总是垂直于其速度和磁感强度构成的平面，具体的指向还与运动电荷自身的电荷量的正负有关，即
\begin{itemize}
    \item 对于正电荷，其磁场力是右手四指由速度转向磁感强度时拇指指向的正方向。
    \item 对于负电荷，其磁场力是右手四指由速度转向磁感强度时拇指指向的负方向。
\end{itemize}
由于洛伦兹力总是垂直于电荷的运动方向，因此，\textbf{洛伦兹力只能改变运动电荷的运动方向}，而并不改变运动电荷的速率和动能，这也相当于是说洛伦兹力不对运动电荷做功。

\subsection{带电粒子在匀强磁场中的运动}
现在我们来研究一种特殊情况，即带电粒子在匀强磁场中的运动。为了便于分析，我们将带电粒子的速度分解为平行于磁感强度和垂直于磁感强度两个分量，随后再将运动合成。
\begin{Figure}{带电粒子在匀强磁场中的运动}
    \inputfigure[scale=1]{Chapter10D_Figure01}
\end{Figure}
如果我们记速度和磁感强度的夹角为$\theta$，平行分量即为$v\sin\theta$，垂直分量即为$v\cos\theta$。

对于平行于磁感强度的速度分量$\vb*{v}\cos\theta$，其并不会导致任何的洛伦兹力，因此在平行于磁感强度的方向上，带电粒子将保持进入磁场前的$v\cos\theta$的速度做匀速直线运动。

对于垂直于磁感强度的速度分量$\vb*{v}\sin\theta$，其会使得带电粒子受到一个垂直于该速度分量的洛伦兹力，而且由于这里的磁场是匀强的，其受到的洛伦兹力的大小并不改变，因此在垂直于磁感强度的方向上，带电粒子将以$v\sin\theta$的速度做匀速圆周运动、如果我们将这两个方向的运动叠加合成，我们会看到，\textbf{带电粒子在匀强磁场中的运动轨迹为一条等速螺旋线}。

\begin{Figure}{带电粒子在匀强磁场中的运动分解}
    \subfigure[主视图]
    {
        \inputfigure[scale=1]{Chapter10D_Figure03}
    }\qquad\qquad
    \subfigure[俯视图]
    {
        \inputfigure[scale=1]{Chapter10D_Figure02}
    }
\end{Figure}

现在我们来具体分析一下这个运动，设带电粒子的质量为$m$，设带电粒子的电荷量为$q$，假设圆周运动的半径为$R$，根据\Refx{公式}{洛伦兹力}和\Refx{公式}{向心加速度}
\begin{Equation}
    F=qv\sin\theta B=\frac{mv^2\sin^2\theta}{R}
\end{Equation}
这样就解得了匀速圆周运动的半径，称为\uwave{回旋半径}
\begin{BoxEquation}[回旋半径]
    R=\frac{mv\sin\theta}{qB}
\end{BoxEquation}
进而我们就可以解得周期
\begin{equation*}
    T=\frac{2\pi R}{v\sin\theta}=\frac{2\pi mv\sin\theta}{qBv\sin\theta}=\frac{2\pi m}{qB}
\end{equation*}
这样就算得了匀速圆周运动的周期，称为\uwave{回旋周期}，其倒数称为\uwave{回旋频率}。
\begin{BoxEquation}[回旋周期]
    T=\frac{2\pi m}{qB}
\end{BoxEquation}
\begin{BoxEquation}[回旋频率]
    \nu=\frac{qB}{2\pi m}
\end{BoxEquation}
\Refx{公式}{回旋周期}指出，带电粒子的回旋周期与其速度无关，其仅取决于磁感强度和带电粒子的质量与电荷量，也就是说，速度的大小和方向只会改变回旋半径和稍后将介绍的回旋螺距，不影响带电粒子的回旋周期，这一结论是磁聚焦和回旋加速器的理论基础。

\Refx{公式}{回旋周期}中我们看到，回旋周期与磁感强度和电荷量成正比，回旋周期与质量成反比，这一结果可以这样理解，较强的磁场意味着电荷更容易被磁场“拽住”，从而回旋的空间就小一些了,另外，更大的电荷量意味着磁场力更强，更大的质量意味着运动状态更难被磁场力改变。有时也会将$m/q$称为带电粒子的\uwave{质荷比}，这样就可以简单的说，\textbf{回旋周期与质荷比和磁感强度成正比}，不过理论中更为常用的是$q/m$即带电粒子的\uwave{荷质比}。

带电粒子在磁场中回旋一周就会前进一定距离，这一距离称为\uwave{回旋螺距}，不难想象，回旋螺距的大小即为带电粒子在平行磁场方向的速度$v\cos\theta$与回旋周期$T$的乘积，即
\begin{BoxEquation}[回旋螺距]
    H=2\pi\frac{mv\cos\theta}{qB}
\end{BoxEquation}
\Refx{公式}{回旋螺距}和\Refx{公式}{回旋半径}中，我们看到，\textbf{回旋半径和回旋螺距均与质荷比和磁感强度成正比}，但两者同时也与速度成正比，并且和速度的方向有关
\begin{itemize}
    \item 速度一定时，速度与磁感强度的夹角越大，回旋半径越大。
    \item 速度一定时，速度与磁感强度的夹角越小，回旋螺距越大。
\end{itemize}
由此可见，某种意义上，速度的方向决定了速度在回旋半径和回旋螺距上的分配。

这里有一个简单的应用，我们知道，虽然我们将电子束假象为一条射线，但实际上电子的角度总归会存在一定的微小差异，这就会使得电子束在光屏上的影像变得模糊，这时我们就可以沿着电子束的方向施加一个匀强磁场，设电子束中的电子的运动速度均为$v$，但与电子束的整体方向存在一个夹角，设电子的运动方向与磁场的夹角为$\theta$，由于这里的$\theta$很小
\begin{equation*}
    v\cos\theta=v\qquad v\sin\theta=v\theta
\end{equation*}
这里我们看到，角度不同的电子在垂直磁场方向的分量$v\theta$不同，但是这些电子平行磁场的分量$v$基本相同，而各个电子的回旋周期又相同，这就意味着，尽管各个电子在磁场中以不同的回旋半径进行运动，但是，在时间上相隔一个回旋周期，在空间上相隔一个回旋螺距，这些电子会重新会聚到同一位置，这样就可以显著的改善电子束的成像，由于这种现象与光束经过透镜后聚焦的现象类似，因此将这种通过磁场会聚带电粒子的方法称为\uwave{磁聚焦}。

下面我们来看一些通过带电粒子在匀强磁场中的运动特性制造的实验设备。

\subsection{回旋加速器}
\uwave{回旋加速器}是原子核物理学中获得高速粒子的一种装置，这种装置的基本原理就是利用回旋周期与速率无关的性质，回旋加速器的核心部分是两个D形盒，它的形状类似于扁圆的金属盒沿直径剖开得到的左右两半，在D形盒之间留有狭缝并存在交变的电场，在D形盒内部由于导体空腔对电场的屏蔽效应，几乎不存在电场。整个装置被置于一个巨大的真空容器中，在垂直D形盒的方向上用电磁铁产生一个均匀磁场，覆盖整个D形盒空间。

当带电粒子最初位于狭缝时，通过电场使其取得一个速度$v_1$并进入右侧的D形盒，在磁场的洛伦兹力作用下以回旋半径$R_1$经半个圆周回到狭缝，这时期望电场恰好反向，从而可以继续加速带电粒子，使其以更大的速度$v_2$在左侧D形盒中以$R_2$为半径再次回旋。
\begin{Figure}{回旋加速器的最初两轮回旋}
    \subfigure[回旋加速器的第一轮加速]
    {
        \inputfigure[scale=0.45]{Chapter10D_Figure04}
    }\qquad
    \subfigure[回旋加速器的第二轮加速]
    {
        \inputfigure[scale=0.45]{Chapter10D_Figure05}
    }
\end{Figure}
根据\Refx{公式}{回旋半径}，我们看到
\begin{equation*}
    R_1=\frac{mv_1}{qB}\qquad R_2=\frac{mv_2}{qB}
\end{equation*}
回旋加速器在设计上的巧妙之处在于，虽然回旋半径随着速度的增大而不断增大，但是回旋周期却并不变化，这也就是说，为了保证带电粒子在每一次经过狭缝时都得到加速，只需要使用一个交流频率为带电粒子的回旋频率的电源就可以了，而不需要人为动态的根据粒子的速度调节电场方向。在交流电场的作用下，带电粒子会在狭缝中被不断加速并取得越来越大的回旋半径，带电粒子最终会接近D形盒的边缘，将其引出后就可以进行相关实验。
\begin{Figure}{回旋加速器}
    \inputfigure[scale=0.7]{Chapter10D_Figure06}
\end{Figure}
回旋加速器能使得带电粒子取得的最大速度，取决于其D形盒半径的大小，因此如果我们需要获得速度很高的带电粒子，就需要建造庞大的回旋加速度器，其半径尺度可能达到千米量级，另外，质荷比较小的带电粒子可以在半径一定的回旋加速器中取得更大的速度。

回旋加速器也存在一定缺陷，这就是当带电粒子的速度接近光速时，由于相对论效应的作用，带电粒子的质量开始增大，带电粒子的质荷比也开始增大，这就会改变带电粒子的回旋频率，这时交流电源的频率与之不相匹配，从而无法进一步加速带电粒子（因为这时带电粒子经过狭缝区域时，有时电场方向与运动方向相同，有时电场方向与运动方向相反）。

如果要进行更进一步的加速，就需要使用原理不同的\uwave{同步加速器}。世界上最大的同步加速器是大型强子对撞机\footnote{简称LHC，即Large Hadron Collider。}，其位于一个圆周达到27公里的环形隧道内，贯穿法国和瑞士边境。

\subsection{速度选择器}
\uwave{速度选择器}是一种可以控制通过的带电粒子速度的设备，它由两块带电平行金属板产生的均匀电场$\vb*{E}$和与之垂直的均匀磁场$\vb*{B}$构成，从粒子源进入速度选择器的带电粒子将会同时受到电场力$\vb*{F}_e=qE$和磁场力$\vb*{F}_m=qvB$的作用，两个力的方向相反，因此
\begin{itemize}
    \item 如果电场力大于磁场力，那么带电粒子将坠向右侧的负极板，无法通过速度选择器。
    \item 如果磁场力大于电场力，那么带电粒子将坠向左侧的正极板，无法通过速度选择器。
\end{itemize}
由此可见，那些可以通过速度选择器的带电粒子必然受到相同的电场力和磁场力。
\begin{Figure}{速度选择器}
    \inputfigure[scale=0.7]{Chapter10D_Figure07}\vspace{-5pt}
\end{Figure}\vspace{-5pt}
由于电场力和磁场力应相等，故
\begin{equation*}
    qE=qvB
\end{equation*}
这就是说，只有速度满足电场和磁场之比的带电粒子才能无偏转的通过速度选择器
\begin{BoxEquation}[速度选择器]
    v=\frac{E}{B}
\end{BoxEquation}
这样我们就可以通过控制电场和磁场的强度之比，来获取我们所需的特定速度的带电粒子。

\subsection{质谱仪}
同种元素的质子数是一定的，但可以具有不同的中子数，因此可以有不同的原子量，这就是所谓的\uwave{同位素}，例如对于氢元素，它具有三种同位素：氕（\ce{_1^1H}）、氘（\ce{_1^2H}）、氚（\ce{_1^3H}）。
\begin{Figure}{氢元素的同位素}
    \subfigure[氕（\ce{_1^1H}）]
    {
        \inputfigure[scale=2]{Chapter10D_Figure08}
    }\qquad
    \subfigure[氘（\ce{_1^2H}）]
    {
        \inputfigure[scale=2]{Chapter10D_Figure09}
    }\qquad
    \subfigure[氚（\ce{_1^3H}）]
    {
        \inputfigure[scale=2]{Chapter10D_Figure10}
    }\qquad
\end{Figure}
同位素的测定可以通过\uwave{质谱仪}实现，首先我们要将待测的同位素混合物电离，使其以离子状态存在，随后这些离子需要经过速度选择器，假设控制的通过速度为$v$，根据\Refx{公式}{回旋半径}，由于经过速度选择器的离子具有相同的速度，因此离子在磁场中的回旋半径就只与离子的质荷比有关了，质荷比越大回旋半径越大，质荷比越小回旋半径越小。
\begin{Figure}{质谱仪}
    \inputfigure[scale=0.7]{Chapter10D_Figure11}
\end{Figure}
如果我们在离子进入磁场回旋半周的位置放置一块感光胶片（或类似传感器），由于不同质荷比的同位素具有不同的回旋半径，因此同位素离子落在胶片上的位置就不同。根据胶片的感光位置就可以推演出各个同位素的质量值，根据感光强弱还可以推出同位素的比例。

如果记质谱仪中的磁场为$\vb*{B}$，而记速度选择器中的磁场和电场为$\vb*{B}_0,\vb*{E_0}$，根据\Refx{公式}{回旋半径}和\Refx{公式}{速度选择器}，经过质谱仪的离子的回旋半径满足以下公式
\begin{BoxEquation}[质谱仪]
    R=\frac{m}{qB}\frac{E_0}{B_0}
\end{BoxEquation}

\subsection{霍尔效应}
实验表明，当我们将一块矩形导体板垂直放置在匀强磁场中时，在矩形的长度方向上通以电流时，在矩形的宽度方向的两个面间会存在电势差，这种现象称为\uwave{霍尔效应}。
\begin{Figure}{霍尔效应}
    \subfigure[正载流子的霍尔效应]
    {
        \inputfigure[scale=0.65]{Chapter10D_Figure12}
    }\qquad
    \subfigure[负载流子的霍尔效应]
    {
        \inputfigure[scale=0.65]{Chapter10D_Figure13}
    }
\end{Figure}
现在我们来研究霍尔效应是如何产生的，我们先考虑载流子为正电荷的情况，正电荷原先依照电流方向在横向做定向运动，但是在磁场$\vb*{B}$的作用下，正电荷会在纵向上发生漂移，这样的结果便是，在导体的1面将富集正电荷，在导体的2面将对应的带负电荷，从而在导体内部形成了一个附加电场$\vb*{E}_H$，称为\uwave{霍尔电场}，其导致的电势差就称为\uwave{霍尔电势差}。

由于霍尔电场力$F_e=qE_H$与磁场力$F_m=qvB$的方向相反，因此霍尔电场的出现会减弱正电荷在纵向上的漂移，随着霍尔电场的增强，最终霍尔电场力和磁场力将达到平衡
\begin{equation*}
    qE_H=qvB
\end{equation*}
这时正电荷也不再发生纵向漂移，因此在导体中就形成了一个稳定的霍尔电场
\begin{BoxEquation}[霍尔电场的速度表示]
    E_H=vB
\end{BoxEquation}
这时的霍尔电场是匀强的，如果记1,2面间的距离为$l$，那么1,2面间的霍尔电势差即为
\begin{BoxEquation}[霍尔电势差的速度表示]
    U_H=vBl
\end{BoxEquation}
这样就解释了霍尔效应产生的原因，不过在通常情况下我们并不知道导体中载流子的运动速度，而我们知道的是导体中通过的电流，因此我们需要将以上两个公式用电流表示。

根据\Refx{公式}{电流的微观表述}
\begin{equation*}
    v=\frac{I}{nqS}
\end{equation*}

这里的$S$是导体的横截面积，我们已经设导体的宽度为$l$，如果记导体的厚度为$d$，那么
\begin{equation*}
    v=\frac{I}{nqld}
\end{equation*}

上式代入\Refx{公式}{霍尔电场的速度表示}和\Refx{公式}{霍尔电势差的速度表示}即得
\begin{Equation}[霍尔效应1]
    E_H=\frac{IB}{nqld}\qquad U_H=\frac{IB}{nqd}
\end{Equation}
上式中的$1/nq$是仅与导体材料有关的系数，因为$n$是载流子的数密度而$q$是载流子的电荷量，因此我们可以将$1/nq$定义为一个关于材料的系数，称为\uwave{霍尔系数}，记作$R_H$
\begin{BoxDefinition}[霍尔系数]
    R_H=\frac{1}{nq}
\end{BoxDefinition}
运用\Refx{定义}{霍尔系数}，式\ref{式:霍尔效应1}就可以表述为
\begin{BoxEquation}[霍尔电场]
    E_H=\frac{IB}{nqld}=\frac{IB}{ld}R_H
\end{BoxEquation}
\begin{BoxEquation}[霍尔电势差]
    U_H=\frac{IB}{nqd}=\frac{IB}{d}R_H
\end{BoxEquation}
由此我们可以看到，当材料的霍尔系数越大，当施加的电流和磁感强度越强，当使用的导体越薄，霍尔效应就越明显，霍尔电势差与导体宽度无关，霍尔电场则与导体宽度成反比。

以上讨论是对于载流子为正电荷的情况进行的，而对于载流子为负电荷的情况，这时电荷的正负和电荷的运动方向都反向了，因此电荷受到的洛伦兹力的方向是不变的，因而电荷仍然向着1面漂移，但由于这里为负电荷，此时在1面与2面间形成的电场会反向，然而电荷受到的电场力的方向仍然是不变的，因而最终电场力也会与磁场力达成平衡，形成霍尔效应。

这里我们发现了一个重要的现象，在过去，我们常常认为正电荷的定向运动与等量负电荷的反向运动是等效的，就电流确实如此，就霍尔效应中两者受到的磁场力和电场力也是如此，但是，两者在宏观上产生的霍尔电场和霍尔电势差却是刚好相反，也就是说
\begin{center}
    霍尔效应可以分辨导体中载流子的类型
\end{center}
换言之，正电荷的定向运动和负电荷的反向运动在电流上等效，但其霍尔效应是不同的。

这里我们可能会有疑问，我们知道导体中的载流子总是带负电的电子，那么为什么我们还要研究载流子带正电的霍尔效应？这是因为我们如果将研究范围从导体扩大到半导体时，我们会发现，不同类型的半导体材料将表现出不同的霍尔效应，这会有两种结果
\begin{itemize}
    \item 霍尔效应表现为负载流子的材料，称为\uwave{电子型半导体}，其载流子是带负电的电子。
    \item 霍尔效应表现为正载流子的材料，称为\uwave{空穴型半导体}，其载流子是带正电的\uwave{空穴}。
\end{itemize}
这里的空穴虽然在微观上仍然是电子的缺失形成的，但我们却必须认为空穴型半导体中是带正电的空穴在运动，而不能视作是带负电的电子在反向运动，因为只有这样才能满足霍尔效应的实验结果，因此空穴在这里并不是一个简单的等效物理模型，而是满足霍尔效应所必须的客观存在。所以可以说，\textbf{载流子的正负性完全是被霍尔效应的结果所定义的}。

\Refx{定义}{霍尔系数}中，正载流子的霍尔系数为正，负载流子的霍尔系数为负，若规定右手四指由磁场方向转向电流方向时，右手拇指指向的为2面，右手拇指背向的为1面，并定义霍尔场强和霍尔电势差的正值方向为由1面指向2面的方向，那么\Refx{公式}{霍尔电场}和\Refx{公式}{霍尔电势差}就可以普遍的适用正载流子和负载流子的情况了。

\Refx{定义}{霍尔系数}指出，霍尔系数与材料的载流子数密度是负相关的，导体中电子的浓度大霍尔效应其实并不不明显，相反，在半导体中载流子的浓度要低的多，在半导体中的霍尔效应也就要显著的多，因此大部分\uwave{霍尔器件}都是用半导体材料制成。

霍尔效应的应用是非常广泛的，除了判定材料的载流子类型，霍尔效应的另外一个重要用途就是可以测定磁场强度，如果已知电流$I$和材料的霍尔系数$R_H$以及材料厚度$d$，且测得霍尔电势差为$U_H$，磁场强度就可以根据\Refx{公式}{霍尔电势差}通过计算得出
\begin{equation*}
    B=\frac{U_Hd}{R_HIB}
\end{equation*}
磁传感器通常就是利用霍尔效应和霍尔器件制成的（这比测定电流受到的力要容易的多）。




\section{载流导线在磁场中的受力}

\subsection{安培力}
在磁场中的载流导线也会受到磁场力的作用，这种力称为\uwave{安培力}。而我们知道载流导线是由大量运动电荷构成，而上一节中我们已经知道运动电荷会受到洛伦兹力，因此宏观上的安培力实际上就是微观上的洛伦兹力集体作用的结果，现在我们就来推导安培力的计算公式。

若我们以$\dif{\vb*{F}}$表述电流源受到的安培力，而以$\vb*{F}$表示每个运动电荷的洛伦兹力
\begin{Equation}
    \dif{\vb*{F}}=\vb*{F}nS\dif{l}=\vb*{F}n\dif{\vb*{V}}=\vb*{F}\dif{n}
\end{Equation}
这样就建立了电流元的安培力$\dif{\vb*{F}}$和运动电荷的洛伦兹力$\vb*{F}$的关系。

根据\Refx{公式}{洛伦兹力}和\Refx{公式}{电流元与运动电荷}
\begin{Equation}
    \dif{\vb*{F}}=\vb*{F}\dif{n}=qv\times\vb*{B}\dif{n}=I\dif{\vb*{l}}\times\vb*{B}
\end{Equation}
即
\begin{BoxEquation}[安培力]
    \dif{\vb*{F}}=I\dif{\vb*{l}}\times\vb*{B}
\end{BoxEquation}
\Refx{公式}{安培力}是对电流元而言的，对于完整的载流导线$L$，应以积分形式表示
\begin{BoxEquation}[安培力的积分形式]
    \vb*{F}=\Ilnt{I\dif{\vb*{l}}\times\vb*{B}}
\end{BoxEquation}
现在我们来考虑一些特殊情况，若载流导线是直导线，而磁场又是匀强磁场时，此时的磁感强度$\vb*{B}$可以作为常量提出积分外，而$\dif{\vb*{l}}$的积分即为由导线一端指向另一端的常矢量$\vb*{l}$
\begin{BoxEquation}[匀强磁场中的载流直导线]
    \vb*{F}=I\vb*{l}\times B
\end{BoxEquation}
值得注意的是，对于匀强磁场中任意形状的载流导线，由于$\dif{\vb*{l}}$的积分的结果总是由导线一端指向另一端的矢量，与导线的具体路径无关（因为$f(x,y)=1$是路径无关的），因此
\begin{BoxPrinciple}[匀强磁场中的安培力性质]
    在均匀磁场中，任意弯曲载流导线的受力等同于从其起点到终点通有相等电流的直导线。

    在均匀磁场中，任意闭合载流导线受到的安培力为零。
\end{BoxPrinciple}
特别注意的是以上结论仅适用于匀强磁场，否则$\vb*{B}$无法作为常矢量提出积分外。

\subsection{平行电流间的安培力}
载流直导线可以在磁场中受到安培力，载流直导线也可以同时产生磁场。

在本小节，我们就来计算一组平行无限长载流直导线中的电流间的相互作用。

由于电流产生的磁场总是沿着电流的右手螺旋方向（且在平行电流的任意直线上的磁场是相等的），而电流受到的安培力是右手四指由电流转向磁场时的拇指指向，因此不难得出\footnote{这里的绘图不是很准确，直导线的磁场是距离反比衰减的，图中以匀强磁场绘之，不过这并没有什么妨碍。}

\begin{Figure}{电流间的相互作用力}
    \subfigure[同向电流相吸]
    {
        \inputfigure[scale=0.6]{Chapter10E_Figure01}
    }\qquad
    \subfigure[异向电流相斥]
    {
        \inputfigure[scale=0.6]{Chapter10E_Figure02}
    }\vspace{-10pt}
\end{Figure}
\Refx{图}{电流间的相互作用力}的结论可以表述如下
\begin{itemize}
    \item 两根电流方向\textbf{相同}的无限长平行电流间的磁场力表现为\textbf{吸引力}（同向相吸）。
    \item 两根电流方向\textbf{不同}的无限长平行电流间的磁场力表现为\textbf{排斥力}（异向相斥）。
\end{itemize}
电流中“同向相吸异向相斥”，电荷间“同性相斥异性相吸”，应注意两者是不同的。

现在我们来进一步的做一些定量计算，根据\Refx{公式}{载流无限长直导线的磁场}
\begin{equation*}
    B_1=\frac{\mu_0I_1}{2\pi d}\qquad
    B_2=\frac{\mu_0I_2}{2\pi d}
\end{equation*}
这里的$\vb*{B}_1,\vb*{B}_2$分别是电流$I_1,I_2$在空间中产生的磁场。

现在我们再来计算载流导线受到的磁场力的大小，虽然由另一直导线产生的磁场并不是匀强的，但是在与之平行的直导线所在直线上的各点处的磁感强度都是相等的，因此\Refx{公式}{匀强磁场中的载流直导线}仍然适用。然而其指出直导线受到的磁场力与它的长度成正比，那么对于这里的无限长直导线，两者间的磁场力就是无穷大，这是没有意义的\footnote{其实这也是合理的，就像密度一定的两根无限长的物体间的万有引力是无穷大一样。}，但是我们仍然可以研究无限长直导线中的某一段在另外一根无限长直导线的磁场中受到的安培力，这将是个有限值，在$I_1,I_2$上任取长度为$l_1,l_2$两端，其受力$\vb*{F}_1,\vb*{F}_2$分别为
\begin{equation*}
    \begin{cases}
        F_1=B_2I_1l_1=\dfrac{\mu_0}{2\pi}\dfrac{I_1I_2}{d}l_1\\[5mm]
        F_2=B_1I_2l_2=\dfrac{\mu_0}{2\pi}\dfrac{I_1I_2}{d}l_2
    \end{cases}
\end{equation*}
由此可见，对于一组平行无限长载流直导线，相同长度的导线总是受到相同的力，无论我们研究的这段导线是具体属于其中哪一根载流直导线，因此上式可以统一表述为
\begin{BoxEquation}[平行载流导线间的作用力]
    F=\frac{\mu_0}{2\pi}\frac{I_1I_2}{d}l
\end{BoxEquation}
\Refx{公式}{平行载流导线间的作用力}是意义是，对于两根载有电流$I_1,I_2$的平行无限长直导线，若两根导线相距为$d$，那么对于其中任何一根导线上长度为$l$的导线段，它将会受到大小为$F$的磁场力，若$I_1,I_2$方向相同则$F$为引力，若$I_1,I_2$方向反向则为斥力。

\Refx{公式}{平行载流导线间的作用力}可以定义国际单位制中的电学基本单位安培
\begin{BoxGeneral}[安培][定义]
    真空中两根相距$1$\si{m}的无限长平行圆直导线中通以等量恒定电流时，

    若每$1$\si{m}长度上的导线受到的作用力为$2\times 10^{-7}$\si{N}，则此时每根导线中的电流为$1$\si{A}。
\end{BoxGeneral}
这里定义了安培后，真空磁导率的数值就被确定了
\begin{Equation}
    \mu_0=4\pi\times 10^{-7}\si{N\cdot A^{-2}}
\end{Equation}
这也就解释了为何真空磁导率的数值是如此整齐了，因为其直接基于一个基本单位的定义。

由于安培的大小被确定后库仑的大小也就确定了，因此根据\Refx{定律}{库仑定律}，如果我们现在按照这个定义取两个$1$库仑的点电荷，使其相距$1$米，那么两者间的库仑力是可以通过实验测定的，这样真空电容率的数值也可以被确定了，而实验表明其值为
\begin{Equation}
    \varepsilon_0=8.85\times 10^{-12}\si{C^2\cdot N^{-1}\cdot m^{-2}}
\end{Equation}
这也就表明，真空磁导率和真空电容率之间是存在耦合关系的，两者不是相互独立的。

这里可以设想，如果我们将\Refx{定义}{安培}中同等条件下的电流改定义为$k$\si{A}，那么上面的点电荷的电荷量也对应的变为$k$\si{C}，如果记此时的真空电容率和真空磁导率为$\varepsilon_0',\mu_0'$
\begin{equation*}
    F_e=\frac{1}{4\pi\varepsilon_0'}\frac{q_1'q_2'}{r^2}\qquad 
    F_m=\frac{\mu_0'}{2\pi}\frac{I_1'I_2'}{d}l
\end{equation*}
而在原先定义下有（注意$F_e,F_m$是不会随着安培和库仑的定义而变的）
\begin{equation*}
    F_e=\frac{1}{4\pi\varepsilon_0}\frac{q_1q_2}{r^2}\qquad 
    F_m=\frac{\mu_0}{2\pi}\frac{I_1I_2}{d}l
\end{equation*}
其中新定义中的$q_1',q_2',I_1',I_2'$是原定义中的$q_1,q_2,I_2,I_2$数值上的$k$倍，因而
\begin{equation*}
    \varepsilon_0'=a^2\varepsilon_0\qquad
    \mu_0'=\frac{\mu_0}{a^2}
\end{equation*}
关注到
\begin{equation*}
    \varepsilon_0'\mu_0'=\varepsilon_0\mu_0
\end{equation*}
这也就是说，当我们修改安培的定义时，尽管真空电容率和真空磁导率的数值会发生变化，但是两者的乘积$\varepsilon_0\mu_0$是一个定值，我们后面还将看到$\varepsilon_0\mu_0$实际与光速$c$有关。


\vspace{-10pt}
\section{静磁场中的磁偶极子}
我们知道，对于匀强磁场中的磁偶极子，也就是载流线圈，根据\Refx{定律}{匀强磁场中的安培力性质}可知其受到的磁场力为零，但是磁偶极子在匀强磁场中仍会受到磁力矩的作用。
\begin{Figure}{静磁场中的磁偶极子}
    \inputfigure[scale=0.6]{Chapter10F_Figure01}\vspace{-10pt}
\end{Figure}\vspace{-5pt}
为了研究的方便，我们先假定这里的磁偶极子是一个矩形载流线圈，且当磁偶极矩的方向确定时，该矩形线圈围需绕磁偶极矩旋转至能使得矩形线圈的两条边垂直于磁场的角度，在这种情况下，两条边上的电流垂直纸面（叉点），两条边上的电流平行纸面（灰色实虚线）。

对于平行纸面的两边上的电流，容易看出两者所受到安培力分别垂直纸面向外和向内，大小相等，方向相反，且通过矩形线圈的中心，既不会产生合力，也不会产生合力矩。

对于垂直直面的两边上的电流，虽然两者的安培力是等大反向的，但是两者的安培力并不是共线的，因此会产生力矩的作用，考虑到两者受力的对称性，我们在计算合力矩时只要考虑其中一条边的电流的力矩并乘以二就可以了，这里我们考虑垂直纸面向上的那根电流。

设垂直纸面的电流的导线长为$l$，设垂直纸面的电流的间距为$d$。

根据\Refx{公式}{匀强磁场中的载流直导线}，垂直纸面向上的电流受到的安培力为
\begin{equation*}
    \vb*{F}=I\vb*{l}\times\vb*{B}
\end{equation*}

这里的$\vb*{l}$的方向是垂直直面向上的电流的方向，设其相对于线圈中心的位矢为$\vb*{d}/2$。这里的矩形线圈受到的力矩，是垂直纸面向上的电流产生的力矩的两倍，根据\Refx{定义}{力矩}
\begin{Equation}[磁偶极子的磁力矩1]
    \vb*{M}=2\qty[\frac{\vb*{d}}{2}\times\qty(I\vb*{l}\times\vb*{B})]=\vb*{d}\times\qty(I\vb*{l}\times\vb*{B})
\end{Equation}
根据矢量三重积的结论，通常情况下矢量叉乘是不满足结合律
\begin{align*}
    \vb*{A}\times(\vb*{B}\times\vb*{C})&=(\vb*{A}\cdot\vb*{C})\vb*{B}-(\vb*{A}\cdot\vb*{B})\vb*{C}\\[4mm]
    (\vb*{A}\times\vb*{B})\times\vb*{C}&=(\vb*{C}\cdot\vb*{A})\vb*{B}-(\vb*{B}\cdot\vb*{C})\vb*{A}
\end{align*}

但如果有$\vb*{A}\cdot\vb*{B}=0$且$\vb*{B}\cdot\vb*{C}=0$，上两式的第二项均为零，此时矢量叉乘满足结合律
\begin{equation*}
    \vb*{A}\times(\vb*{B}\times\vb*{C})=
    (\vb*{A}\times\vb*{B})\times\vb*{C}\qquad(\vb*{A}\perp\vb*{B}~~\vb*{B}\perp\vb*{C})
\end{equation*}

而这里的$\vb*{d},\vb*{l}$和$\vb*{l},\vb*{B}$是垂直的，因此式\ref{式:磁偶极子的磁力矩1}可以化为
\begin{Equation}[磁偶极子的磁力矩2]
    \vb*{M}=(I\vb*{d}\times\vb*{l})\times\vb*{B}=(IS\vb*{e}_n)\times\vb*{B}
\end{Equation}
上式中的$S$为线圈面积$S=dl$，而$\vb*{e}_n$恰好是磁偶极矩方向的单位矢量。

根据\Refx{定义}{磁偶极矩}
\begin{BoxEquation}[磁偶极子的磁力矩]
    \vb*{M}=\vb*{m}\times\vb*{B}
\end{BoxEquation}
而对于任意形状的平面线圈构成的磁偶极子，例如我们最为熟悉的圆线圈，我们可以假象其为无数窄矩形线圈的组合\footnote{由于面积上的相等，平面线圈的磁偶极矩是这些载有相同电流的窄矩形线圈的磁偶极矩的和。}，在两个窄矩形线圈的公共边上会有相反流向的电流，其磁场效应必然相会抵消，就像这些公共边不存在一样，而实际上这些公共边确实是不存在的，因此这种设想并不破坏原有的物理图像，因此\Refx{公式}{磁偶极子的磁力矩}总是适用的。

\Refx{公式}{磁偶极子的磁力矩}指出，磁偶极子在匀强磁场中所受\uwave{磁力矩}的方向，是右手四指由磁偶极矩转向磁感强度时右手拇指的方向，而其旋转方向即这里右手四指的转向。

\Refx{公式}{磁偶极子的磁力矩}中，如果我们记$\vb*{m},\vb*{B}$的夹角为$\theta$，那么当$\theta=\pi/2$时磁力矩最大，而当$\theta=0$或$\theta=\pi$时磁力矩为零，但前者是稳定平衡，而后者是不稳定平衡。
\begin{Figure}{磁偶极子的平衡}
    \subfigure[稳定平衡]
    {
        \inputfigure[scale=0.5]{Chapter10F_Figure02}
    }\qquad
    \subfigure[不稳定平衡]
    {
        \inputfigure[scale=0.5]{Chapter10F_Figure03}
    }
\end{Figure}
因此，磁偶极子在磁力矩的作用下，总是会偏转至磁偶极矩与磁感强度一致的方向。
